% file bismon/doc/bismon-chariot-doc.tex; please read both
% bismon/README.md and bismon/dependencies/README-dependencies.md
% before processing this LaTeX document using the "make doc" command
% from bismon/ source directory.
\documentclass[11pt,a4paper,svgnames]{article}
\usepackage[T1]{fontenc}
\usepackage[utf8x]{inputenc}
\usepackage{alltt}
% https://tex.stackexchange.com/a/342804/42406
% IMPORTANT! see also: https://tex.stackexchange.com/a/500963/42406
%\usepackage{textcomp}
\usepackage{moreverb}
\usepackage{fancyvrb}
\usepackage{fancyhdr}
\usepackage{fancybox}
% libertine, see https://tex.stackexchange.com/a/9868/42406
%\usepackage{libertine}
%\usepackage{epsfig}
\usepackage{graphicx}
\usepackage{float}
\usepackage{xcolor}
\usepackage{moreverb}
\usepackage{multirow}
\usepackage{boxedminipage}
\usepackage[square]{natbib}
% https://tex.stackexchange.com/a/16992/42406
\usepackage{mathabx}
%\usepackage{charter}
%\usepackage{inconsolata}
\usepackage{hevea}
\usepackage{listings}
\usepackage{relsize}
\usepackage{verbatimbox}
\usepackage{filecontents}
\usepackage{catchfile}
\usepackage{lastpage}
\usepackage{stmaryrd}
\usepackage{ucs}
%\usepackage{stix}
\usepackage{newunicodechar}
% bigfoot enables \verb in footnotes
\usepackage{bigfoot}
\usepackage{makeidx}
\usepackage{times}
\usepackage[a4paper, margin=2cm]{geometry}
\usepackage{hyperref}

\newcommand{\bmemail}[1]{{\href{mailto:#1}{\texttt{\textbf{#1}}}}}
\newcommand{\bmurl}[1]{{\href{#1}{\texttt{\relsize{-1}{\textbf{#1}}}}}}


% see also http://www.sascha-frank.com/Arrow/latex-arrows.html
% and http://tug.ctan.org/info/symbols/comprehensive/symbols-a4.pdf
% and https://ctan.math.illinois.edu/macros/latex/contrib/newunicodechar/newunicodechar.pdf
%%%% keep in order
%U+21A6 RIGHTWARDS ARROW FROM BAR
\newunicodechar{↦}{$\mapsto$}
%U+21B3 DOWNWARDS ARROW WITH TIP RIGHTWARDS
\newunicodechar{↳}{\rotatebox[origin=c]{180}{$\Lsh$}}
%U+2208 ELEMENT OF
\newunicodechar{∈}{$\in$}
% U+00AB LEFT-POINTING DOUBLE ANGLE QUOTATION MARK
\newunicodechar{«}{\guillemotleft}
% U+00BB RIGHT-POINTING DOUBLE ANGLE QUOTATION MARK
\newunicodechar{»}{\guillemotright}
% U+00B1 PLUS-MINUS SIGN
\newunicodechar{±}{$\pm$}
% U+00B5 MICRO SIGN
\newunicodechar{µ}{$\mu$}


%see https://tex.stackexchange.com/a/51656/42406 
\lstdefinelanguage
   [x64]{Assembler}     % add a "x64" dialect of Assembler
   [x86masm]{Assembler} % based on the "x86masm" dialect
   % with these extra keywords:
   {morekeywords={CDQE,CQO,CMPSQ,CMPXCHG16B,JRCXZ,LODSQ,MOVSXD, %
                  POPFQ,PUSHFQ,SCASQ,STOSQ,IRETQ,RDTSCP,SWAPGS, %
                  rax,rdx,rcx,rbx,rsi,rdi,rsp,rbp, %
                  r8,r8d,r8w,r8b,r9,r9d,r9w,r9b, %
                  r10,r10d,r10w,r10b,r11,r11d,r11w,r11b, %
                  r12,r12d,r12w,r12b,r13,r13d,r13w,r13b, %
                  r14,r14d,r14w,r14b,r15,r15d,r15w,r15b}} % etc.

%% https://latex.org/forum/viewtopic.php?t=3970
%\makeatletter
%\newcommand*\bmexpandableinput[1]{\@@input#1 }
%\makeatother

 % see https://tex.stackexchange.com/a/51349/42406
\hypersetup{
  colorlinks   = true, %Colours links instead of ugly boxes
  urlcolor     = NavyBlue, %Colour for external hyperlinks
  linkcolor    = DarkGreen, %Colour of internal links
  citecolor   = DarkMagenta, %Colour of citations
  frenchlinks = true,
}
% generated files
\input{generated/git-commit.tex}
\input{generated/timestamp.tex}



% LaTeX specific macros that HeVeA should ignore
%% file bismon-latex.tex
% LaTeX specific macros that HeVeA should ignore

\newcommand{\bmincludewidthgraphics}[4]{\includegraphics[width=#1]{generated/#2.#3}}



\date{\bmdocdate}

%-\title{{\color{orange}{D1.3v1}}~Specialized Static Analysis tools for more secure and safer IoT software development.
%-  \\
%-  {\large{\emph{Bismon}~\thanks{The \texttt{bismon} name is temporary and
%-    is likely to change later. See also the
%-    \bmurl{http://github.com/bstarynk/bismon} repository; this documentation was generated with its \texttt{build-bismon-doc.sh} script}~ documentation}}}
%-
%-\author{Basile \textsc{Starynkevitch}\thanks{Email
%-    \bmemail{basile@starynkevitch.net} or
%-    \bmemail{basile.starynkevitch@cea.fr}}, %
%-  \\ {\small{CEA, LIST,
%-      France}}}
%-
% the index works again in commit d7070ccc581cf951e2a22bfed8 but we dont know why :-(
\makeindex

\begin{document}

\bibliographystyle{abbrvnat}


\pagestyle{fancy}
\fancyhf{}
\rhead{\textsc{Chariot} D1.3~\textsuperscript{v2-}}
\lhead{\textmd{Specialized Static Analysis tools for  for more secure and safer \\
  IoT software development} [draft] \hspace{3cm} \emph{(Bismon)}}
\renewcommand{\footrulewidth}{0.4pt}
% https://tex.stackexchange.com/a/317531/42406
\futurelet\TMPfootrule\def\footrule{{\color{orange}\TMPfootrule}}
\futurelet\TMPheadrule\def\headrule{{\color{orange}\TMPheadrule}}
\fancyfoot[L]{{\raisebox{0.0cm}[1pt][1pt]{\color{LightSlateGrey}{{\relsize{+1}{DRAFT}}}~\relsize{-1.5}{\texttt{\bmgitcommit} on \emph{\bmgitdate}}}}}
\fancyfoot[R]{{Page \thepage ~ {\relsize{-2.5}{of \pageref{LastPage}}}}}
\begin{titlepage}
%\maketitle

  \thispagestyle{empty}

  \begin{center}
    
    %   \bmincludewidthgraphics{width=194pt}{CHARIOT-logo-fig}{eps}{eps}
    \includegraphics[width=194pt]{images/CHARIOT-logo}

    {\large \textbf{C}ognitive \textbf{H}eterogeneous  \textbf{Ar}chitecture for
Industrial~\textbf{IoT}}
    
    {\color{orange}\Huge {D1.3~\textsuperscript{v2-}~Specialized Static Analysis tools for more secure and safer IoT software development.}\\
    {\Large \emph{Bismon} documentation.}}

  
\begin{relsize}{-0.1}
\subsection*{Document Summary Information}
  \begin{tabular}{|l|l|l|l|}
    \hline 
  \textbf{Grant Agreement} & 780075 & \textbf{Acronym} & \textsc{Chariot} \\
  \hline 
    {\textbf{Full title}} & %
    \multicolumn{3}{l|}{\textbf{\relsize{+0.5}{C}}ognitive \textbf{\relsize{+0.5}{H}}eterogeneous \textbf{\relsize{+0.5}{A}}rchitecture for Industrial {\relsize{+0.5}{\textbf{IoT}}}}  \\
    \hline
    \textbf{Start date} & 01/01/2018 & \textbf{Duration} & 36 months \\
    \hline
  \textbf{Project URL} & \multicolumn{3}{l|}{{\large \href{http://www.chariotproject.eu}{\textbf{\texttt{www.chariotproject.eu}}}}} \\
    \hline
  \textbf{Deliverable} & \multicolumn{3}{p{9.5cm}|}{{\textbf{D1.3} Specialized Static Analysis tools for more secure and safer IoT software development.}} \\
    \hline
  \textbf{Workpackage} & \multicolumn{3}{l|}{\textbf{WP1}} \\
    \hline
  \textbf{Dissemination level} & \multicolumn{3}{l|}{\textbf{Public report} {\large {[DRAFT]}}} \\
    \hline
  \textbf{Contractual due date} & {31/06/2020} (M30) & \textbf{Actual submission date} & - \\
    \hline
  \textbf{Nature} & R {\small (report)} & \textbf{dissemination level} & PU  {\small (public)} \\
    \hline
 \textbf{Lead beneficiary} & \multicolumn{3}{l|}{CEA, LIST} \\
 \hline
 \textbf{Responsible author} & \multicolumn{3}{l|}{%
   \begin{minipage}{8cm}\smallskip 
     Basile \textsc{Starynkevitch} (CEA, LIST) ~ {\relsize{-1.5}{DILS/LSL}} \\
     {\relsize{-1.5}{from its \textsc{Software Security Laboratory}}} \\
     {\relsize{-1}{\bmemail{basile.starynkevitch@cea.fr}}} \\
     {\relsize{-1}{\bmurl{http://starynkevitch.net/Basile/}}} \\
     \smallskip
   \end{minipage}}
\\
  \hline
  \textbf{Contributions from} &\multicolumn{3}{l|}{ - }
  \\
    \hline
  \textbf{Revision} &\multicolumn{3}{p{9cm}|}{git commit \texttt{\bmgitcommit} on \textit{\bmgitdate}} \\
    \hline 
  {\emph{{\LaTeX}-generated}} &\multicolumn{3}{p{9cm}|}{{\scriptsize\textit{\bmdoctimestamp}}} \\
    \hline
  {{\relsize{-1}{\emph{draft downloadable from}}}} &\multicolumn{3}{p{12cm}|}{{\relsize{-0.5}{\bmurl{http://starynkevitch.net/Basile/bismon-chariot-doc.pdf}}}} \\
    \hline
\end{tabular}
\end{relsize}

\bigskip

\vspace{3cm}

\bigskip

 \begin{tabular}{c l}
   \hline
       \raisebox{-0.3cm}{{\bmincludewidthgraphics{0.08\textwidth}{Flag-of-Europe-fig}{eps}{svg}}}
       &
       \begin{minipage}{0.7\textwidth}\smallskip
         \begin{relsize}{-0.2}
           This \textsc{Chariot} project has received funding from the European Union’s Horizon 2020 research and innovation programme under the Grant Agreement No 780075.
         \end{relsize}\\
         \end{minipage}
       \\
   \hline
 \end{tabular}

  \end{center}

  \newpage


  \begin{center}
      \Ovalbox{
        \begin{minipage}{0.85\textwidth}
          \medskip

    \begin{relsize}{-1.5}
      The astute and busy reader unrelated to the \textsc{Chariot}
      H2020 \index{h2020} project is invited to skip directly to
      section \S \ref{subsec:audience}, since the text written
      immediately below is targetted for evaluators or partners of
      \textsc{Chariot} and its understanding may require access to
      some confidential material.
    \end{relsize}
          
    \medskip
    
    \end{minipage}
  }
  \end{center}
  
  
  \subsection*{Revision history (including peer reviewing \& quality control)}
  
  \begin{tabular}{|p{0.15\textwidth}|p{0.15\textwidth}|p{0.15\textwidth}|l|l|}
    \hline
    \textbf{Version} & \textbf{Issue date} & \textbf{\% Complete} & \textbf{Changes} & \textbf{Contributor(s)} \\
    \hline
    % D1.3 v1 & 30/11/2018 & 100\% & Version 1 & B. Starynkevitch \\
    \multicolumn{5}{|c|}{\parbox[c][1cm][c]{0.5\textwidth}{\Large \textbf{\emph{Not Applicable}}} - see explanation below} \\
    \hline
  \end{tabular}

  \medskip
  
%- {\large {\color{red}{@@For CHARIOT consortium: I (Basile Starynkevitch) could give a few \texttt{git commit}-s but not all of them. If you require them to be listed, then which \emph{precise}
%-       \texttt{git commit}-s could be considered as having been
%-       peer-reviewed and quality-controlled? Please list them
%-       extensively and send an email to me, \href{mailto:basile.starynkevitch@cea.fr}{\texttt{basile.starynkevitch@cea.fr}}. Thanks!}}}
%-
%- \medskip
  
This partly generated document and the \emph{Bismon} software itself
are co-developed in an agile and incremental manner, and have exactly
{\bmgitnumbercommits} \href{http://git-scm.com/}{\texttt{git}} commits
on \textit{\bmdoctimestamp}. For more, please see
\bmurl{https://github.com/bstarynk/bismon/commits/} for details. These
commits are too many and too fine grained to be considered as
``versions''. Given the agile and continuous workflow, it is
unreasonable, and practically impossible, to identify any formalized
versions.

\begin{relsize}{-0.5}
This document is co-developed with the \emph{Bismon} software itself,
it was typeset using \LaTeX~on Linux and contains some \emph{generated}
documentation~\footnote{The generated parts are clearly identified as
  such, and are extracted from the \emph{Bismon} system.}, mixed with
hand-written text. During development of \texttt{bismon}, the amount
of generated documentation will grow.  The entire history of
\emph{Bismon} (both the software -including its persistent store- and
this document) is available on
\bmurl{https://github.com/bstarynk/bismon/commits} and has, for this
document of commit id \texttt{\bmgitcommit} (done on \emph{\bmgitdate}) generated on
\textit{\bmdocdate}, exactly {\bmgitnumbercommits} commits (or
elementary changes). Since changes on any file in the \texttt{git}
repository can affect this document, no ``version'' is identifiable.
\end{relsize}
\bigskip

For convenience to the reader, here are the last~\footnote{Obtained by the
  \texttt{git log --name-status -3} command running in \emph{bismon}
  top source directory.} three \texttt{git commit}-s:

\begin{relsize}{-2}
  \input{generated/000-lastgitcommits.tex}
\end{relsize}

There is no notion of any identifiable ``version'' in \emph{bismon}, so
also in this D1.3~\textsuperscript{v2-} deliverable. The work is incremental and the
development is agile.

\bigskip

\subsection*{disclaimer}

\begin{center}
  \fbox{\parbox{0.92\textwidth}{ \small The content of the publication
      herein is the sole responsibility of the publishers and it does
      not necessarily represent the views expressed by the European
      Commission or its services.

      While the information contained in
      the documents is believed to be accurate, the authors(s) or any
      other participant in the \textsc{Chariot} consortium make no
      warranty of any kind with regard to this material including, but
      not limited to the implied warranties of merchantability and
      fitness for a particular purpose.

      Neither the \textsc{Chariot}
      Consortium nor any of its members, their officers, employees or
      agents shall be responsible or liable in negligence or otherwise
      howsoever in respect of any inaccuracy or omission herein.
      
      Without derogating from the generality of the foregoing neither
      the \textsc{Chariot} Consortium nor any of its members, their
      officers, employees or agents shall be liable for any direct or
      indirect or consequential loss or damage caused by or arising
      from any information advice or inaccuracy or omission herein.  }
  }
\end{center}

\smallskip



 \subsection*{copyright message}

 \begin{center}
   \Ovalbox{
     \parbox{0.95\textwidth}{
     \medskip
     
    
       Copyright © 2018 - 2019 CEA {\relsize{-0.2}{(Commissariat à
           l'énergie atomique et aux énergies alternatives)}}.

       \bigskip 
         %
     This deliverable contains original unpublished work except where
     clearly indicated otherwise. Acknowledgement of previously
     published material and of the work of others has been made
     through appropriate citation, quotation or
     both. Reproduction is
     authorised provided the source is acknowledged  {\relsize{-0.5}{(see last page for licensing details)}}.
    

     \medskip
     }
     }
 \end{center}

 \smallskip
 

  \subsection*{Executive Summary}

  This D1.3\textsuperscript{v1} \textsc{Chariot} deliverable is
  {\color{red}{(was)}} a first \textbf{draft} -and \emph{preliminary}-
  version (at M12) of D1.3\textsuperscript{v2} scheduled at M30 on
  ``Specialized Static Analysis tools for more secure and safer IoT
  software development''.

But {\color{red}{\textbf{incremental work has begun on
      D1.3\textsuperscript{v2} in early 2019}}}, hence the
D1.3\textsuperscript{v2-} in the header.

It describes the \textsc{Chariot} vision on static source code (mostly
of C and C++ code for IoT firmware and application) analysis. It
proposes a \emph{simple} static analysis \emph{framework} leveraging
on the powerful recent \emph{GCC} [cross-]compiler. A \emph{persistent
  monitor} (tentatively named \texttt{bismon}) is being designed and
implemented as a GPLv3+ free software for Linux, using
meta-programming techniques (leveraging on experience gained in the
former GCC MELT project) to \emph{generate} GCC plugins, and able to
keep some intermediate results (of compilation or static analysis)
during the entire life of the IoT project, and giving to the IoT
developers (thru a web interface) a whole-program view of the source
code (as digested by the \emph{GCC} cross-compiler) and of its static
analysis properties. That framework is configurable and scriptable by
static analysis experts, hence permitting different IoT projects to
address various concerns, while keeping the usual IoT development
workflow (running as usual their \emph{GCC} cross-compiler on Linux,
with extra plugin-related compilation flags).  The deliverable has
been structured starting from the identification of the software and
tool users and the document expected audience as well as the vision on
specialized source code analysis towards more secure and safer IoT
software development. The report then describes its strong alignment
to adding capabilities to \emph{GCC} as well as the driving principles of the
tools. Data and their persistence character are also described
including mutable and non-mutable values/types while persistence is
considered to start by loading some previous persisted state, usually
dumping its current state before termination and loading the next
state on the next load-up. The framework for static code analysis is
also defined as part of the \emph{GCC} compilation process. The described
work is also analysed in terms of contributing to other free software
projects.

\medskip

\bigskip

\subsection*{notice}

\begin{quote}
\begin{relsize}{-0.5}
This work \index{h2020@\emph{\textsc{h2020}}} is funded\footnote{The development work of the
  \href{http://github.com/bstarynk/bismon}{\texttt{bismon}} free
  software prototype is also \emph{partly} funded -from 2019 to 2021-
  by the \href{http://decoder-project.eu/}{\textsc{Decoder}} H2020
  project, under its Grant Agreement 824231.} (from start of 2018 to
end of 2020) thru the \emph{\textsc{Chariot}} project (see its web
site on \bmurl{http://chariotproject.eu/}) which has received funding
from the European Union’s Horizon 2020 research and innovation
programme under the Grant Agreement No 780075.
\end{relsize}
\end{quote}
\medskip



\hspace{2cm}

\begin{center}
%{\small git commit: \texttt{\bmgitcommit}}
%
%{\small generated: \textit{\bmdoctimestamp}}
%
%\hspace{1cm}

\begin{tabular}{cc}
  \bmincludewidthgraphics{72pt}{CHARIOT-logo-img}{png}{png} %
  & \bmincludewidthgraphics{64pt}{Flag-of-Europe-fig}{eps}{svg}
\end{tabular}


\hspace{1cm}

\hrule

\bigskip

\textbf{Deliverable created during the \href{http://chariotproject.eu/}{\textsc{Chariot}} project.}
\end{center}
\end{titlepage}





\newpage

\tableofcontents

\newpage
\listoffigures

\medskip

\listoftables

\newpage

\rfoot{Page \thepage ~ {\relsize{-2}{of \pageref{LastPage}}}}
\newpage


\subsection*{Glossary of terms and abbreviations used}

\smallskip

\begin{tabular}{|c|p{0.7\textwidth}|}
  \hline
  \begin{minipage}{0.15\textwidth}

    \smallskip
    
    \textbf{Abbreviation / Term}

    \smallskip
    
  \end{minipage} & \textbf{Description} \\
  \hline \emph{binutils} \index{binutils@\emph{binutils}} & GNU free software package containing assembler \index{as@\texttt{as}} \texttt{as}, linker \texttt{ld} \index{ld@\texttt{ld}} and other utilities to operate on object files or binary executables, etc... \bmurl{https://www.gnu.org/software/binutils/} \\
  \hline \index{bismon@\texttt{bismon}} \texttt{bismon} & the free software framework and persistent monitor described here ; source repository on \bmurl{http://github.com/bstarynk/bismon/} \\
  \hline \textsc{Clang} \index{clang@\textsc{Clang}} & The Clang open-source project provides a language front-end and tooling infrastructure for languages in the C language family (C, C++, Objective C/C++, OpenCL, CUDA, and RenderScript) for the LLVM project \bmurl{http://clang.llvm.org/} \\
  \hline \textsc{Frama-C} \index{Frama-C@\textsc{Frama-C}} &  a free sofware extensible platform for analysis of C software \bmurl{http://frama-c.com/} \\
  \hline FSF & Free Software Foundation \bmurl{http://fsf.org/} \\
  \hline GCC \index{GCC} & Gnu Compiler Collection \bmurl{http://gcc.gnu.org/} \\
  \hline \textsc{Gcc Melt} \index{MELT@\textsc{Melt}} \index{GCC MELT@\textsc{Gcc Melt}} & was a (GPLv3+-licensed) GCC plugin and framework providing a DSL \index{DSL} to ease GCC extensions ; it is archived on \bmurl{http://starynkevitch.net/Basile/gcc-melt/}\\
  \hline \emph{Generic} \index{generic@\emph{Generic}} & language-independent abstract syntax tree (internal representation) in GCC \\
  \hline Gimple \index{gimple@\emph{Gimple}} & middle-end internal representation in GCC \\
  \hline GPL\index{GPL} & Gnu General Public Licence (a copylefted free software license) \bmurl{https://www.gnu.org/licenses/gpl.html} \\
  \hline IoT \index{IoT} & Internet of Things \\
  \hline \texttt{libonion} \index{libonion@\texttt{libonion}} & an HTTP server library \bmurl{https://www.coralbits.com/libonion/} \\
  \hline LLVM \index{LLVM} & The LLVM Project is an open-source collection of modular and reusable compiler and toolchain technologies \bmurl{http://www.llvm.org/} \\
  \hline MELT \index{MELT} & the Lisp-like domain specific language used in GCC MELT \\
  \hline Persistence \index{persistence} & From Wikipedia~: ``In computer science, \emph{persistence} refers to the characteristic of state that outlives the process that created it. This is achieved in practice by storing the state as data in computer data storage''. \bmurl{https://en.wikipedia.org/wiki/Persistence\_(computer\_science)} \\
  \hline RTL & (register transfer language) back-end internal representation in GCC \\
  \hline static code analysis \index{static code analysis} \index{code analysis} \index{analysis} & (or static program \index{program!analysis} analysis) ``is the analysis of computer software that is performed without actually executing programs, in contrast with dynamic analysis, which is analysis performed on programs while they are executing.'' (from Wikipedia: \bmurl{https://en.wikipedia.org/wiki/Static\_program\_analysis}). In this D1.3\textsuperscript{v1} report, it means static source code analysis, in practice analysis of C or C++ code for IoT fed to the \emph{GCC} compiler. \\
  \hline SSA \index{SSA} & Static Single Assignment (in GCC, a variant of Gimple) \\
  \hline
\end{tabular}

%%%%%%%%%%%%%%%%%%%%%%%%%%%%%%%%%%%%%%%%%%%%%%%%%%%%%%%%%%%%%%%%
%%%% the chapters
\newpage
% file intro-bm.tex
\section{Introduction}

\subsection{About \textit{Bismon}}

\textit{Bismon}\footnote{\texttt{bismon} is a \textbf{temporary name}
  and could be changed, once we find a better name for it} is a free
software static source code whole-program analysis framework (whose
initial domain will be \emph{Internet of Things} (or
\index{IoT}{IoT})\footnote{IoT is viewed as the first application
  domain of \textit{Bismon}, but it is hoped that most of
  \textit{Bismon} could be reused and later extended for other
  domains}. It is designed to work with the \textit{Gcc} compiler (see
\bmurl{gcc.gnu.org}) on a Linux workstation. \textit{Bismon} is the
successor of \textit{GCC MELT} \footnote{The \textit{GCC MELT} web
  pages used to be on \texttt{gcc-melt.org} -a DNS domain relinquished
  in april 2018- and are archived on
  \bmurl{starynkevitch.net/Basile/gcc-melt}} \cite{Starynkevitch2007Multistage, starynkevitch-DSL2011}.

\textit{Bismon} is \textbf{work in progress}, and many things
described here are not fully implemented in 2018.




\newpage

%% file datapersist-bm.tex, which is \input from bismon-doc.tex
\section{Data and its persistence in \emph{Bismon}}

\subsection{data processed in \emph{Bismon}}

The \emph{Bismon monitor} handles various kinds of
\index{data}{data}. A lot of data is \index{immutable}{immutable} (its
content cannot change once the data has been created, for example
strings). But \index{object}{\textbf{objects}} are of course mutable
and can be modified after creation. Since \emph{Bismon} is
multi-threaded and its agenda is running \emph{several} worker threads
in parallel, these mutable objects contain a \index{mutex}{mutex} for
locking purposes.

So the \emph{Bismon monitor} handle values: they can be immutable, or
objects (and objects are the only kind of mutable data).

\subsubsection{Immutable values}

They include

\begin{itemize}

\item \textbf{tagged integer} (of 63 bits). The \index{integer}{integer} won't change
  (and integer values don't require extra space to keep that integer).
  
\item UTF-8 encoded \textbf{string} (the bytes inside such
  \index{string}{strings} don't change).

  \item \textbf{tuple} of objects, that is an ordered (but immutable)
    sequence of object pointers (the size or content of a tuple don't
    change). A given object could appear in several positions in a
    \index{tuple}{tuple}.

    \item \textbf{set} of objects, organized internally as a sorted
      array of object pointers. A given object can occur only once in
      a set, and membership (of an object inside a set) is tested
      dichotomically in logarithmic time. Of course, the size and
      content of a \index{set}{set} don't change with time.

      \item \textbf{node}. A \index{node}{node} has an object
        connective, and a (possibly empty, but fixed) sequence of
        \index{son}{sons} (sons are themselves values, so can
        themselves be integers, strings, tuples, sets, sub-nodes). The
        connective, size and sons of a node don't change with
        time. Since a node is immutable and knows all its sons at
        creation time, circularity (e.g. a node having itself as some
        grand-son) inside it is impossible, and the node has a finite
        depth.
        
\end{itemize}

The \index{nil}{\textbf{nil}} value is generally not considered as a
value, but as the absence of some value. We might (later) add other
kind of values (perhaps vectors of integers, of doubles, bignums ...),
but they should all be immutable. It is likely that we prefer complex
or weird data to sit inside objects.

Tuples and nodes could contain nil, but sets cannot. A node connective
is a genuine object (so cannot be nil), but a node could have nil
sons.

Sets and tuples are sometimes both considered as
\index{sequence}{\textbf{sequences}} and share some common operations.

The immutable values are somehow lightweight. Most of them (strings,
sets, tuples, nodes) internally keep some hash-code (e.g. to
accelerate equality tests on them, or accessing hash tables having
values as keys). The memory overhead for values is therefore small (a
few extra words at most, to keep GC-data type and mark, size hash).

The size of values (byte length of strings, number of objects in
tuples or sets, number of sons in nodes) can in principle reach
$2^{31} - 1$ but is generally much smaller (often less than a few
dozens) and could be 0.

Mutable values outside of objects (and their payload, see below)
cannot exist.

Values (even references to objects, e.g. inside sequences or nodes)
are represented as a machine pointer and fit in a 64 bits word. When
its least significant bit is 1, it is a tagged integer.

Values, including objects, are comparable so sortable. For nodes,
sets, tuples we use a lexicographical order.

\subsubsection{Mutable objects}

\index{object}{Objects} are the only kind of \index{mutable}{mutable}
values, and are somehow heavy (at least severy dozens of machine words
in memory). They can be accessed nearly simultaneously by several
worker threads running different tasklets, so they need a locking
mechanism and contain a (recursive) mutex\footnote{Created with
  \href{http://man7.org/linux/man-pages/man3/pthread\_mutex\_lock.3p.html}{\texttt{pthread\_mutex\_lock}(3)}
  with \texttt{PTHREAD\_MUTEX\_RECURSIVE} attribute} (so in reality
only one thread is using them at a given instant).

\medskip

\emph{Conceptually}, objects contain the following data:

\begin{itemize}
  \item a constant \emph{unique} serial id (of about 128 bits), called
    the \index{objid}{\textbf{objid}}, randomly generated at object
    creation time and never changed after. In many occasions, that
    objid is printed as 24 characters (two glued blocks of 12
    characters, the first being an underscore \texttt{\_}, the second
    being a digit, the 10 others being alphanumerical with significant
    case) such as \texttt{\_4ggW2XwfXdp\_1XRSvOvZqTC}~\footnote{That
      objid \texttt{\_4ggW2XwfXdp\_1XRSvOvZqTC} is for the predefined
      object \emph{\texttt{the\_system}}, and corresponds to the two
      64 bits numbers 3577488711679049683, 1649751471969277032 i.e. to
      128 bits hexadecimal 0x31a5cb0767997fd316e5183916681468} or
    \texttt{\_0xbmmxnN8E8\_0ZuEqJmqMNH}. It is expected that objid
    \index{collision}{collisions} never occur in practice, e.g. that
    even thousands of \emph{Bismon monitor} processes (running on many
    distant computers) would in fact never generate the same objid. In
    other words, our objids are as unique as \index{UUID}{UUID}s (from
    \href{https://tools.ietf.org/html/rfc4122}{RFC 4122}) are (but are
    displayed differently, without hyphens). The concrete textual
    ``syntax'' for objid-s is carefully chosen to be compatible and
    friendly with identifiers in C, C++, Java, Ocaml, etc. The
    \emph{Bismon} runtime maintains a large array of hashtables and
    mutexes to be able to quickly find the object pointer of a given
    objid (if such an object exists in memory). The objid is used to
    compare (and sort) objects. The (32 bits, non-zero) hash code of
    an object is obtained from its objid (but it is cached in the
    object's memory, for performance reasons).

    \item the recursive \textbf{mutex lock} of that object\footnote{We
      have considered using a pthread \texttt{rwlock} instead of a
      \texttt{mutex}, but that would probably be more heavy and
      perhaps slower, but could be experimented in the future.}. So
      locking (or unlocking) an object really means using that lock on
      pthread mutex operations~\footnote{So accessing without the
        protection of that lock being hold, any data inside an object,
        other than its constant objid, its class, its routine pointer,
        is forbidden and considered as \emph{undefined behavior}}.

  \item a \index{space}{\textbf{space}} number fitting in a single
    byte. The space 0 is for \index{transient}{\emph{transient}}
    objects that are not persisted to disk. The space 1 is for
    \index{predefined}{\emph{predefined}} objects (there are about 60
    of them in Q3 of 2018), which are conceptually created before the
    start of \emph{Bismon monitor} processes and are permanently
    available, even at initial load time of the persistent
    store. Those predefined objects are dumped in file
    \texttt{store1.bmon}, the objects of space 2 (conventionally
    called the \emph{global} space) are dumped and persisted in file
    \texttt{store2.bmon}, etc...

    \item the \index{mtime}{\textbf{mtime}} of an object holds its
      modification time, with a millisecond granularity, since the
      Unix Epoch. \index{touch}{Touching} an object is updating its
      \emph{mtime} to the current time.

    \item the (mutable!) \index{class}{\textbf{class}} of an object is
      an \index{atomic}{\emph{atomic}}~\footnote{``atomic'' in the C
        or C++ memory sense; so a pointer declared \texttt{\_Atomic}
        in C or \texttt{std::atomic} in C++, supposing that they are
        the same and interoperable.} pointer to an object (usually
      another one) describing its class, as understood by
      \emph{Bismon}. It is allowed to change
      dynamically~\footnote{Changing classes is permitted with
        reasonable bounds: the class of all classes should remain the
        \emph{\texttt{class}} predefined object, but of course such an
        invariant cannot be proved.} the class of any object. Classes
      describe the behavior (i.e. the dictionnary of ``methods''), not
      the content (i.e. the ``attributes'') of objects and enable
      single-inheritance (every class has one super-class).

    \item the \index{attribute}{\emph{attributes}} of an object are
      organized as an hash-table associating attribute or key objects
      to arbitrary non-nil values. An \textbf{attribute} is an
      arbitrary object, and its value is arbitrary (but cannot be
      nil).

    \item the \index{component}{\emph{components}} of an object are
      organized as a vector (whose size can change, grow, or shrink)
      of values. A \textbf{component} inside an object is a value
      (possibly nil).

      \item objects may contain one \textbf{routine} pointer (or nil),
        described by \begin{enumerate}
          
        \item the \index{routine address}{\emph{routine address}}
          inside an object is a function or \index{routine}{routine}
          pointer (in the C sense, possibly and often nil). The
          signature of that function is described by the routine
          signature~\footnote{Perhaps all our routines will keep the
            same signature, and then it would not need to be
            explicitly stored.}

        \item the \index{routine signature}{\emph{routine signature}}
          is (when the routine address is non-nil) describing the
          signature of the routine address above.
      \end{enumerate}

      Notice that routine address and signature can only change when
      a new module~\footnote{The generated C code of modules also
        contains an array of constant objids, ana another array of
        routine objids.} is loaded (or at initial persistent state load
      time), and that can happen only when the agenda is
      inactive. Conceptually they are mostly constant (and do not
      require any locking).
        
      Most (in 2018, all) routines have the same C signature
      \texttt{objrout\_sigBM} corresponding to the predefined object
      \texttt{\emph{function\_sig}}. For an object of objid
      \texttt{\emph{oid}} of that signature its routine address
      corresponds to the C name \texttt{crout\emph{oid}\_BM}. For
      instance, to initialize (at load time) the object of oid
      \texttt{\_9CG8SKNs6Ql\_4PiHd8cnydn} the initial loader (or the
      module loader) would \texttt{dlsym} the
      \texttt{crout\_09Hug4WGnPK\_7PpZby8pz84\_BM} C function name.


      \item objects may also have some (nearly arbitrary)
        \index{payload}{\textbf{payload}}. That is a pointer (possibly
        nil) to some client data owned by the object; the payload is
        usually not a value but something else. The garbage collector
        should know all the payload types. In 2018 the following payloads are possible:

        \begin{enumerate}
        \item mutable \index{vector}{\emph{vector}} of values (like for components)
        \item mutable \index{associative table}{\emph{associative
            table}} associating objects to non-nil values (like for
          attributes)
          \item mutable \index{hash set}{\emph{hash set}} of objects.
        \end{enumerate}
\end{itemize}


\newpage

% file staticanalys-bm.tex, which is \input from bismon-doc.tex
\section{Static analysis of source code in \emph{Bismon}}
\label{sec:staticanalys}

Static analysis involves a \emph{generated} GCC plugin (whose C++ code
is generated by the \emph{bismon} persistent monitor) which
communicates with the monitor and sends to it some digested form of
the analyzed C or C++ code. Some translation-unit specific processing
can happen in that GCC plugin, but the whole program aspects of the
static code analysis should obviously be done inside the monitor, and
requires -and justifies- its persistence. The complexity and
non-stability of \emph{GCC} internal representations justify some
semi-automatic approach in extracting them (see \S\ref{subsec:analygcc}
below).

%The rest of this chapter will be written in the final D1.3\textsuperscript{v2} version.

%-A significant part of this chapter should be generated (like \emph{GCC
%-  MELT} generated its documentation, see
%-\cite{Starynkevitch-GCCMELTweb}) from the persistent state of
%-\emph{Bismon}. Perhaps this chapter should be put after the ``using
%-Bismon'' chapter (\S\ref{sec:using}).

\subsection{static analysis of \emph{GCC} code}
\label{subsec:analygcc}

The \emph{GCC} compiler has a complex (and ill-defined,
under-documented and evolving, so unstable) application programming
interface (API) which can be used by plugins. So \emph{Bismon} needs
to analyze the various \emph{GCC} plugin related \emph{header files}
to extract important information about that API, so to be later able
to generate \emph{GCC} plugin code. Such an extraction (inspired by
the approach inside \emph{Clasp}, which does similar things with the
help of \emph{Clang}, see \cite{Schafmeister:2015:CLASP} for details)
needs not to care about the \emph{Gimple} instructions, but only about
the abstract syntax tree in \emph{Tree} and \emph{Generic} forms (see
\cite{gcc-internals} \S11) to retrieve the full description of
\emph{GCC}.

This approach of extracting semi-automatically~\footnote{We are well
  aware that some work still needs to be done manually, in particular
  giving the really useful subpart of the \emph{GCC} API.} the GCC API
(of parsing GCC header files with some simple GCC plugin) is motivated
by past GCC MELT experience (where every feature of the GCC API had to
be \emph{explicitly} and manually described in MELT language; these
descriptions took a lot of time to be written and had to be manually
maintained; however, most of them could in theory be extracted
automatically from GCC headers).

A bootstrapping and incremental approach, in several ``steps'', is
worthwhile (and possible because of persistence): we will first
extract some very simple information from GCC header files, and use
them to improve the next extraction from the same GCC header
files. The \emph{slow} evolution~\footnote{GCC internals are
  \emph{slowly} evolving, because GCC itself is huge: its
  ``navigation'' is as slow as that of a supertanker which needs hours
  to turn and change directions. So for \emph{social} reasons the GCC
  community is changing the API slowly, but there is no promise of
  stability.}  of GCC API is practically relevant (most of the API of
\texttt{gcc-8.3} should stay in the next \texttt{gcc-9.0} version).

Descriptive data related to the API of a particular version of GCC
will thus stay persistently in the \emph{Bismon} monitor, but should
be updated at each release of \emph{GCC}. We care mostly about API
related to optimization passes, \emph{GENERIC}, \emph{Gimple},
\emph{SSA} and \emph{Optimized-Gimple}. We probably don't need to go
at the \emph{RTL} level. The version
\href{https://gcc.gnu.org/gcc-10/}{10 of \textsc{Gcc}} (released in
May 2020) incorporates several
\href{https://gcc.gnu.org/onlinedocs/gcc/Static-Analyzer-Options.html}{static
  analysis options}, that are activated with the \texttt{-fanalyzer}
option to \texttt{g++} or \texttt{gcc}. The version
\href{http://lists.llvm.org/pipermail/llvm-announce/2020-March/000087.html}{10
  of \textsc{Clang}} (released in March 2020) contains an improved
\href{https://clang-analyzer.llvm.org/}{Clang static analyzer} and
\href{https://clang.llvm.org/extra/clang-tidy/}{\texttt{clang-tidy}}
``linter'' like tool, check both coding styles and portability
related conventions. Both compilers are open source and available
\footnote{Both recent \textsc{Gcc} and \textsc{Clang} are buildable as
cross compiler for major 32 bits or 64 bits architectures such as
PowerPC, x86, x86-64, or ARM, provided one download their source
code.} on a \textsc{Linux} desktop and both should be of interest for
advanced \emph{IoT} software developers coding in C or in C++ and
capable of using the command line. \textsc{Gcc} analysis features be
configured thru appropriate \texttt{\#pragma}-s, and \textsc{Clang}
analysis features are changeable by conventional comments such as
\texttt{// NOLINT(google-explicit-constructor, google-runtime-int)}.

%Additional content of this \S\ref{subsec:analygcc} will be written
%for the final D1.3\textsuperscript{v2}.


\subsection{static analysis of IoT firmware or application code}
\label{subsec:analysiot}

Once the API of the current version of \emph{GCC} is known to the
persistent monitor, we can generate the C++ code of \emph{GCC} plugins
for cross-compilers used by IoT developers.

A first static analysis, useful to IoT developers, will be related to
whole-program detection of \emph{stack overflow} \index{stack
  overflow} (see also \cite{Payer:2018:MSV}). By the way, such an
analysis is currently not doable by \emph{Frama-C}, because it don't
know the size of each call frame. However, \emph{GCC} is already
computing that size (see the \texttt{-fstack-usage} option which dumps
the size of the call frame of each function, and the
\texttt{-Wframe-larger-than=\emph{bytesize}} option), and we simply
need to extract and keep it. We also need to get a good approximation
of the \emph{control flow graph\index{control flow graph}}. For that
we need to extract basic blocks and just \texttt{GIMPLE\_CALL}
\emph{Gimple} statements (ignoring other kinds of \emph{Gimple}
statements). Of course, indirect calls (thru function pointers, which
are infrequently used in most IoT code) are harder to handle (and
could require interaction with the IoT developer using our monitor, to
annotate them).

A proof-of-concept GCC plugin for GCC 8 (and 9) to take advantage of
existing internal GCC passes to compute some upper approximation of
the call stack size has been developped. That hand-written GCC plugin,
coded in file \texttt{gcc8plugin-demo-chariot-2019Q2.cc} of about a
thousand lines of C++, communicate with the \texttt{bismon} monitor
using some REST HTTP protocol with ad-hoc HTTP \texttt{POST} requests
having a JSON payload, in some \textsc{Chariot} specific JSON
format. The \texttt{bismon} monitor should display that diagnostic in
a Web browser tab. It could also use the \emph{language server
  protocol}~\footnote{See \bmurl{https://langserver.org/} for more.}
which is, in 2019, understood by most free software source code
editors running on Linux, including \texttt{emacs}, or \texttt{vim},
or \texttt{VSCode}. It might even later use the new
\emph{Sarif}\footnote{See
  \bmurl{http://docs.oasis-open.org/sarif/sarif/v2.0/csprd01/sarif-v2.0-csprd01.html}
  for more.}  protocol, designed for communication between static
source code analyser.

Notice that according to
\href{https://www.zdnet.com/article/chrome-70-of-all-security-bugs-are-memory-safety-issues/}{this
  webpage}, nearly 70\% of security bugs affecting the Chrome web
browser (by Google) are related to memory management issues in C++.
It is expected that junior European software developers of
non-critical IoT systems would experiment a similar bug distribution
in their IoT code. A long-term approach could be the costly training
of IoT software engineers to switch to programming languages with
better memory management, such as \href{https://golang.org/}{Go} or
\href{https://www.rust-lang.org/}{Rust}. However, rewriting an entire
IoT code base is too costly, and mixing several programming
languages\footnote{See
\href{https://softwareengineering.stackexchange.com/questions/370135/why-are-multiple-programming-languages-used-in-the-development-of-one-product-or}{this}
for a discussion of why is that interesting. Notice that recent
\href{http://gcc.gnu.org/}{GCC} compilers share some common internal
representations between several language front-ends.} in the same
software product can be worthwhile but requires some rare and
qualified labor. \href{https://gcc.gnu.org/gcc-10}{GCC 10} has a new
\href{https://gcc.gnu.org/onlinedocs/gcc/Static-Analyzer-Options.html}{static
  analysis} framework (with its \texttt{-fanalyzer} compiler option)
and powerful warning options\footnote{It is helpful to pass
\texttt{-Wall -Wextra} to the \texttt{gcc} or \texttt{g++} compiler,
usually thru some
\href{https://en.wikipedia.org/wiki/Build_automation}{build
  automation} tool such as
\href{http://ninja-build.org/}{\texttt{ninja}}}. The
\href{https://clang-analyzer.llvm.org/}{\textsc{Clang} static
  analyzer} could be useful. Some coding rules (such as
\cite{Holzmann:2006:power-of-10} or \index{Misra@\textsc{Misra C}}
\href{https://www.misra.org.uk/}{\textsc{Misra C}}) are available, but
the \href{https://en.wikipedia.org/wiki/Rice's_theorem}{Rice's
  theorem} \index{Rice's theorem} forbids the possibility of a sound
and complete static analyzer: false alarms cannot be avoided, and code \index{code review}
reviews by senior programmers is still necessary.

We probably would also take as an example the analysis of some
\href{http://mqtt.org/}{MQTT library}. The insight is to trust some
existing MQTT implementation~\footnote{Our purpose is not to prove the
  correctness of a given MQTT implementation, which would require a
  formal methods approach à la \textsc{Vessedia}, but to help the
  developer using and trusting it, by checking some specific coding
  rules.}, and to help \emph{junior} developers in using it, by
checking simple coding rules relevant to MQTT.

An interesting \textsc{Chariot}-compatible approach could be to use
\emph{topological data analysis} \index{topological data analysis}
\index{data analysis!topological} (cf
\cite{Chazal:2017:topodatanalys}) techniques, combined with some
machine learning (cf \cite{flach:2012:machine-learning}) and big data
\index{big data} / data mining \index{data mining} (cf
\cite{wu:2013:big-data-mining, clarke:2016:big-data-risks,
  zuboff:2015:big-other, helbing:2019:big-data-democracy}) approaches,
on some of the several directed graphs (notably the \emph{control flow
graph}, the \emph{call graph}, the \emph{dependency graph} for
example) of the whole analyzed program. Reputable free software
libraries\footnote{See \textsc{TensorFlow} on
\bmurl{https://www.tensorflow.org/}, and \textsc{Gudhi} on
\bmurl{http://gudhi.gforge.inria.fr/}, and many other similar
libraries.} are available on Linux. In principle, such an approach
might be used in \texttt{bismon} for a \emph{semi-automatic} detection
of \index{code smell} \index{smell!code} \emph{code smells}. Sadly,
the lack of allocated human resources, and the strong focus (see
\cite{Heder:2017:TRL}) \index{TRL} on high
TRL\footnote{\emph{Technical Readiness Level} and the
\href{https://en.wikipedia.org/wiki/Technology\_readiness_level}{TRL
  wikipage} for more.} results, forbids even trying such an
interesting approach in \textsc{Chariot}, taking into account that
industrial corporations are not even dreaming of it. However, these
approaches might be tried in some other projects, perhaps
\index{decoder@\textsc{Decoder}}
\href{https://www.decoder-project.eu/}{\textsc{Decoder}}.

\medskip
\subsection{static analysis related to pointers and addresses}
\label{subsec:analysptr}

Pointers \index{pointer} are an important part of the C11 and C++11
language specifications (see \cite{C11:std, CplusPlus11:std})}, but
  are difficult to understand. They need several chapters\footnote{A
    novice programmer should be explained that after \texttt{int
      tab[4]; int* p = \&tab+1;} both \texttt{tab[2]} and
    \index{alias!pointer} \texttt{p[1]} are pointer aliases, but
    \texttt{sizeof(tab)} is not \texttt{sizeof(p-1)} even if
    \texttt{tab == p-1}.}  in C or C++ textbooks such as
  \cite{gusted:t2019:modern, Stroustrup:2014:CplusPlus}. Practically
  \index{null@\texttt{NULL} in C}
  \index{nullptr@\texttt{\textbf{nullptr}} in C++} speaking, a pointer
  is an address, with \texttt{NULL} (in C) or
  \texttt{\textbf{nullptr}} (in C++) having a special and
  distinguished meaning: it is never the same as the result of the
  \index{address} \emph{address-of} operator (unary \texttt{\&} prefix
  operator).

 From the IoT or firmware developer's point of view, pointers -viewed
 as addresses- may behave strangely in practice, and differently from
 the language specifications: in theory, the \texttt{NULL} pointer
 might not sit at address 0. In practice, IoT or firmware developers
 do know that (on most implementations) it \emph{is} at address
 0. Dereferencing the \texttt{NULL} pointer is the prototypical
 example of \emph{undefined behavior}, yet some firmware
 code\footnote{A typical example would be some \textsc{Bios} or
   \textsc{Uefi} firmware or operating system kernel on most PC
   desktop motherboards, see
   \href{https://osdev.org/}{\texttt{osdev.org}} and
   \href{http://tinyvga.com/}{\texttt{tinyvga.com}} for more} may do
 that with good reasons. For example, some
 \href{https://en.wikipedia.org/wiki/AVR_microcontrollers}{\textsc{Avr}
   microcontrollers} \index{avr@\textsc{Avr} microcontroller} (used in
 \index{arduino@\textsc{Arduino} platform}
 \href{https://arduino.cc/}{\textsc{Arduino} platforms} ``the working
 registers are mapped in as the first 32 memory addresses'' (from the
 \href{https://en.wikipedia.org/wiki/AVR_microcontrollers}{AVR
   microcontrollers} wikipage). The
 \href{https://en.wikipedia.org/wiki/MIPS_architecture}{\textsc{Mips}
   architecture} (used by some \textsc{Chariot} partners) has an
 \index{mips@\textsc{Mips} architecture} instruction set without
 proper I/O ports (see
 \href{https://www2.cs.duke.edu/courses/fall13/compsci250/MIPS32\_QRC.pdf}{this}),
 so input/output happens by accessing some dedicated memory locations:
 from the IoT programmer's point of view, physical I/O happens by
 writing into documented memory locations. The opensource
 \href{https://riscv.org/}{\textsc{Risc-V} architecture}
 \index{riscv@\textsc{Risc-V} architecture}, used in IoT (see
 \cite{lee:2020:miot, waterman:2016:riscv-design})) also has
 memory-based input/output, but admit extensions with separate
 input/output ports.

 Most IoT used operating system kernels (see
 \cite{ArpaciDusseau14-Book}, \index{Linux@\textsc{Linux} operating
   system} \index{kernel!operating system} \index{operating system}
 \index{FreeBSD@\textsc{FreeBSD} operating system}
 \index{FreeRTOS@\textsc{FreeRTOS} kernel}
 \href{https://osdev.org}{\texttt{osdev.org}}) such as the
 \textsc{Linux} \href{http://kernel.org/}{kernel} (see also
 \href{https://kernelnewbies.org}{\texttt{kernelnewbies.org}} and
 \href{https://stackoverflow.com/}{\texttt{stackoverflow.com}}),
 \href{https://www.freebsd.org/}{\textsc{FreeBSD}} and
 \href{https://freertos.org/}{\textsc{FreeRTOS}} are managing
 \href{https://en.wikipedia.org/wiki/Process_(computing)}{processes}
 having each their own
 \href{https://en.wikipedia.org/wiki/Virtual_address_space}{virtual
   address space}and provide the
 \href{https://en.wikipedia.org/wiki/Virtual_memory}{virtual memory}
 abstraction, with some
 \href{https://en.wikipedia.org/wiki/File_system}{file systems} to
 application code.

 \medskip

 In most C or C++ source programs, either in some \emph{freestanding
 environment} (like operating system kernels, low-end embedded
 software without any \texttt{main}, e.g. cheap
 \index{arduino@\textsc{Arduino} platform} \textsc{Arduino}-like
 devices) or in a \emph{hosted environment} (so using the standard C
 or C++ libraries and started thru their \texttt{main} function), for
 example some web server running under \textsc{Linux} on something
 similar to a \index{raspberrypi@\textsc{RaspBerryPi} system}
 \textsc{RaspBerryPi} board), dynamically allocated memory is very
 important in practice. See also \cite{karvinen:2014:make} and
 \href{https://www.raspberrypi.org/}{\texttt{raspberrypi.org}} and
 \href{https://www.arduino.cc/}{\texttt{arduino.cc}} websites.

 In hosted environments, dynamic memory allocation often uses
 functions like
 \href{https://man7.org/linux/man-pages/man3/malloc.3.html}{\texttt{malloc}}
 or its friends \texttt{calloc} with \texttt{realloc} and
 \texttt{free}. On some operating systems, it could happen thru
 lower-level system calls such as
 \href{https://man7.org/linux/man-pages/man2/mmap.2.html}{\texttt{mmap}}
 or
 \href{https://man7.org/linux/man-pages/man2/sbrk.2.html}{\texttt{sbrk}}. Deallocation
 would use other system calls such as
 \href{https://man7.org/linux/man-pages/man2/munmap.2.html}{\texttt{munmap}}.

 Many freestanding environments provide more or less equivalent memory
 allocation facilities. For instance,
 \index{FreeRTOS@\textsc{FreeRTOS} kernel} the \textsc{FreeRTOS}
 kernel has in some cases a \texttt{pvPortMalloc} function with a
 behavior close to \texttt{malloc}, and the
 \index{Linux@\textsc{Linux} kernel} the \textsc{Linux} kernel
 \href{http://www.jikos.cz/jikos/Kmalloc_Internals.html}{uses
   \texttt{kmalloc}}, with other deallocating functions close in
 behavior to \texttt{free}. But on \index{arduino@\textsc{Arduino}
   platform} \textsc{Arduino} devices using \texttt{malloc} is indeed
 \href{https://arduino.stackexchange.com/q/682/12068}{discouraged}.

 Dynamic memory allocation is often used but relevant for two
 important concerns: \index{buffer overflow} \index{call stack}
 \index{stack overflow}
 \href{https://en.wikipedia.org/wiki/Buffer_overflow}{buffer
   overflows} and \href{https://en.wikipedia.org/wiki/Call_stack}{call
   stack} management, including avoidance of
 \href{https://en.wikipedia.org/wiki/Stack_overflow}{stack
   overflow}. Call stacks are needed for any multi-threading or
 multi-tasking approach. Junior developers may easily write code that
 blows them up (e.g. with a too naive recursion in C, or by having
 huge automatic variables). A rule of thumb is to avoid needing, in
 hosted environments, call stacks above a megabyte, or call frames
 bigger than a few kilobytes. With recent \texttt{gcc} compilers,
 \index{wstackusage@\texttt{-Wstack-usage} warning option} passing
 something like \texttt{-Wstack-usage=2048} to the \texttt{gcc}
 compilation command is practically helpful.

 Of course, call stacks are a scare resource, and even on Linux it is
 generally unreasonable to have many thousands of them in the same
 process. This explains why \index{thread} \emph{thread}-s are
 expensive, and why the \index{pthread@\textsc{Pthread} library}
 \href{https://computing.llnl.gov/tutorials/pthreads/}{\textsc{Pthread}}
 function \texttt{pthread\_create} (or C++11 \texttt{std::thread}-s)
 should be used with caution. See also
 \href{https://man7.org/linux/man-pages/man7/pthreads.7.html}{\texttt{pthreads(7)}},
   \href{https://man7.org/linux/man-pages/man7/nptl.7.html}{\texttt{nptl(7)}}
     and the source code of \index{glibc@\textsc{Glibc} library}
     \index{musl@\textsc{Musl libc} library}
     \href{https://www.gnu.org/software/libc/}{GNU \textsc{Glibc}} and
     of \href{https://musl.libc.org/}{\texttt{musl libc}}.

 Depending on the development efforts (so costs) available and on the
 criticity of the IoT software,
 \href{https://en.wikipedia.org/wiki/Out_of_memory}{out of memory
   conditions} \index{out of memory} are handled differently. A lot of
 C code incorrectly assumes that \index{malloc@\texttt{malloc}}
 \texttt{malloc} always\footnote{See
 \href{https://stackoverflow.com/a/8460584/841108}{this
   joke-implementation of \texttt{malloc}}, conforming to the letter
 but not the spirit of the C standard.} succeed. In practice, this
 assumption is generally\footnote{On Linux, be aware of
 \href{http://www.etalabs.net/overcommit.html}{memory overcommit}
 which could \href{https://unix.stackexchange.com/q/441364/50557}{be
   disliked}.}  correct, so for non-critical IoT software\footnote{As
 a consumer, I guess that most consumer rented Internet boxes -
 generally running some Linux - do have some system-wide memory leaks,
 and that observation could explain why rebooting them weekly improve
 the user experience.} it does have some economical sense.


 Most non-critical IoT software (e.g. inside consumer devices, perhaps
 \index{raspberrypi@\textsc{RaspBerryPi} system} \textsc{RaspBerryPi}
 based, or wifi routers based upon \index{openwrt@\textsc{OpenWrt}
   router} \href{https://openwrt.org/}{\textsc{OpenWrt} router}
 software) won't care about and won't even handle (or even try to
 cleverly detect and report) rare error conditions or failures
 \index{failure!of standard C functions} such as :

 \begin{itemize}

 \item rare out of memory conditions \index{memory} (e.g. \texttt{malloc} failures)

 \item file system \index{file system} errors (e.g. \texttt{fopen}
   failure on some \index{configuration file} \index{file system}
   \index{etc@\texttt{/etc/} configuration directory} configuration
   file under \texttt{/etc/}), including I/O errors \index{I/O error}
   \index{fscanf@\texttt{fscanf} function} \index{fread@\texttt{fread}
     function} (e.g. \texttt{fscanf} or \texttt{fread} failures) due
   \index{USB storage} to hardware failure (like malfunctioning USB
   storage keys).

     \item floating point \index{floating point!precision} precision
       issues. \index{precision} See the \index{stance@\textsc{Stance}
         project} \href{http://www.stance-project.eu/}{\textsc{Stance}
         European project} (grant 317753), the
       \href{https://floating-point-gui.de/}{\texttt{floating-point-gui.de}}
       website, and \cite{Kiss:2015:combining, goubault:2011:static}.

     \item losing \index{time} time. See the
       \index{time7@\texttt{time(7)}}
       \href{https://man7.org/linux/man-pages/man7/time.7.html}{\texttt{time(7)}}
       \texttt{man} page for an interesting overview. A battery
       powered \index{clock!real-time}
       \href{https://en.wikipedia.org/wiki/Real-time_clock}{real-time
         clock} chip is cheap, but in most IoT consumer devices
       (typically routers or printers) it is not useful enough, since
       they are often connected to some distant
       \index{ntp@\textsc{NTP} - network time protocol}
       \href{https://en.wikipedia.org/wiki/Ntpd}{\textsc{NTPd}
         service}

       \item synchronization bugs related to multi-threading or
         concurrent systems, e.g. with \textsc{Posix threads}, See
         \cite{goubault:2005:practical, sangiorgi:2003:pi,
           david:2013:everything,
           guerraoui:2018:concurrent-systems}. Concurrency is also a
         concern of client-server architectures, including
         \href{https://en.wikipedia.org/wiki/Web_service}{Web
           services} \index{web service} using
         \href{https://en.wikipedia.org/wiki/Hypertext_Transfer_Protocol}{HTTP}
         \index{http@\emph{HTTP} (web protocol)},
         \href{https://en.wikipedia.org/wiki/Simple_Mail_Transfer_Protocol}{SMTP}
         \index{smtp@\emph{SMTP} (mail sending protocol)},
         \href{https://en.wikipedia.org/wiki/Internet_Message_Access_Protocol}{IMAP}
         \index{imap@\emph{IMAP} (email retrieval protocol)}
         \index{email} (perhaps mixed with
         \href{https://en.wikipedia.org/wiki/Transport_Layer_Security}{TLS}
         \index{tls@\emph{TLS} cryptographic protocol}
         \index{cryptographic protocol}).
         
 \end{itemize}

 Pragmatically good reasons to avoid testing such error conditions
 \index{error condition} include: 1. cost of additional development
 efforts for extra C code which is difficult to test in a reproducible
 way \index{test} \index{fault injection} \index{reproducing!faults}
 without fault injections; 2. scarse code memory space \index{code}
 \index{memory!code} (e.g. in some read-only memory) \index{read-only
   memory} on cheap devices.

 Please notice that above error conditions are related: for example, a
 loss of time may corrupt an entire file system,
 \index{corruption!file system} and timing information (or meta-data)
 \index{timing} \index{meta-data} is quite often kept and needed in
 small \href{http://sqlite.org/}{\textsc{Sqlite} databases}.
 \index{sqlite@\textsc{Sqlite} database library}. The remote backup
 \index{database} \index{remote backup} \index{backup!remote} of such
 small or large databases (see \cite{Date:2005:Database-in-Depth,
   kornacker:2015:impala}) could happen on a periodical basis (see
 \href{https://man7.org/linux/man-pages/man5/crontab.5.html}{\texttt{crontab(5)}}
 and \index{crontab@\texttt{crontab(5)}}
 \index{rsync@\texttt{rsync(1)}}
 \href{https://man7.org/linux/man-pages/man1/rsync.1.html}{\texttt{rsync(1)}}
 for
 more). \href{https://softwareengineering.stackexchange.com/a/332071/40065}{Large}
 relational databases are practically needed for most machine learning
 \index{machine learning} algorithms (see
 \cite{flach:2012:machine-learning}), which are usually concurrent
 \index{concurrency} programs. So a non-critical machine-learning
 distributed IoT system would probably be made of cheap IoT devices
 communicating with some powerful, semi-centralized, database
 \index{weather prediction} \index{meteorological!computing grids}
 \index{computing grids} servers. Current meteorological computing
 grids for
 \href{https://en.wikipedia.org/wiki/Numerical_weather_prediction}{numerical
   weather prediction} operated by national meteorological services
 such as \href{http://meteofrance.com/}{Météo-France} require such
 continent-wide IoT networks, and so does \index{forecast!avalanche}
 \index{avalanche} avalanche detection and forecasting in European
 \index{smart cities} \index{city!smart} \index{grid!smart}
 \index{power!grid}
 mountains. \href{https://en.wikipedia.org/wiki/Smart_city}{Smart
   cities} networks and
 \href{https://en.wikipedia.org/wiki/Smart_grid}{smart grids} (see
 \cite{belarbi:2004:vehicle, mclaren:2015:sharing,
   delons:2008:pirandello, bakken:2014:smart-grids} and
 \index{\textsc{Grid4eu} project} the
 \href{https://www.enedis.fr/grid4eu}{\textsc{Grid4Eu} project}) also
 need a large grid of many communicating IoT devices.

 In \emph{some} even critical IoT devices (e.g. medical devices
 \index{medical device} such \index{Covid-19@\textsc{Covid 19}
   illness} as a \href{https://github.com/Recovid/}{Covid-19 breathing
   ventilator}), losing power or time for a few seconds is acceptable,
 as long as there is some hardware alarm (e.g. electronic bell)
 informing professional users (e.g. medical nurses).

 Most non-critical IoT devices could \emph{sometimes} access by the
 network a remote database \index{database} \index{remote!database},
 such as some \href{https://www.postgresql.org/}{\textsc{PostGreSQL}
   server}, or some \href{https://www.mongodb.com/}{\textsc{MongoDB}
   document database}. Losing the connection to such a remote database
 is generally affordable, if the connection loss don't last too
 long. On the other hand, a posteriori adding database checkpoint
 \index{checkpointing!application} \index{scientific code} \index{oil
   industry} \index{automotive crash} \index{super-compuiters}
 \index{code refactoring} \index{refactoring!code}
 \index{code!scientific} \index{code!refactoring} facilities to some
 existing large long-running scientific code (think of oil industry
 simulation software, or digital twins for automotive crash
 \index{simulation software} \index{twins!digital} simulations, both
 running for months of super computer time) requires some significant
 development efforts or code refactoring. Of course European particle
 accelerators such as
 \href{https://home.cern/science/accelerators/}{CERN \textsc{Elena} or
   \textsc{Lhc} installations} are deploying wide networks of
 heterogenous IoT devices, but probably can afford partial failure of
 some of them. Cyber-attacks on smart grids
 (e.g. \cite{lee:2016:analysis-cyberattack}) could justify an increase
 of European research funding on long-term mixed static analysis and
 dynamic machine learning based techniques, but require funding of
 projects with a time span above three years and involving tight
 cooperation thru open source projects (see \cite{Brooks:1995:MM,
   Lerner-Tirole:2000:economics-open-source,
   Tirole:2018:eco-bien-commun, hashem:2015:rise}). See also the related
 \index{softwareheritage@\textsc{SoftwareHeritage} project}
 \href{https://www.softwareheritage.org/}{\textsc{SoftwareHeritage}}
 project. See
 \href{https://en.wikipedia.org/wiki/Brooks's_law}{Brook's law} and
 observation -in 1975- that ``while it takes one woman nine months to
 make one baby, "nine women can't make a baby in one month"'' (to be
 scaled in 2020 by a factor of 10x and generalized to software
 developers of both genders, given the complexity of current software
 intensive systems).

 Notice that exact static prediction of call stack depth is in
 practice impossible, because of
 \href{https://en.wikipedia.org/wiki/Rice's_theorem}{Rice's theorem}
 \index{rice@Rice's theorem} -and usage in C code of the
 \href{https://man7.org/linux/man-pages/man3/alloca.3.html}{\texttt{alloca(3)}}
 stack allocation primitive, or of
 \href{https://en.wikipedia.org/wiki/Variable-length_array}{variable-length
   arrays} - and since compilers may put automatic variables in
 \index{asm@\texttt{asm} keyword} processor registers. For C or C++
 compilers handling the \texttt{asm} keyword (in recent \textsc{Gcc}
 compilers, it provides even
 \href{https://gcc.gnu.org/onlinedocs/gcc/Using-Assembly-Language-with-C.html}{powerful
   language extensions} to C or C++), register allocation becomes a
 nightmare for the compiler writer, since the \texttt{register}
 \index{register@\texttt{register} keyword} keyword tend to become
 obsolete or ill-defined \index{volatile@\texttt{volatile} keyword}
 (like the \texttt{volatile} one). So it is strongly tied to the
 \href{https://en.wikipedia.org/wiki/Register_allocation}{register
   allocation} issue, which is a difficult sub-problem inside most
 optimizing compilers (see \cite{eisl:2016:traceregalloc,
   Chaitin:1981:regalloc, Aho:2006:DragonBook}).  \index{register
   allocation} \index{alloca@\texttt{alloca} stack-allocation
   primitive}. In recent C++ language dialects, the \texttt{auto} and
 \index{auto@\texttt{auto} keyword} \index{decltype@\texttt{decltype}
   keyword} \texttt{decltype} keywords are enabling some \emph{type
 inference} \index{type inference} (see \cite{pierce:2002:types,
   Stroustrup:2014:CplusPlus, Stroustrup:2020:thriving,
   CplusPlus11:std}). In practice, industrial compilers don't even try
 to find the best register allocation possible, but just a good enough
 one (otherwise compilation time could be prohibitive). Be aware that
 \href{https://en.wikipedia.org/wiki/Flexible_array_member}{flexible
   array members} are also a C language feature, which, combined with
 \index{cast!pointer} \index{union@\texttt{union}} pointer casts (or
 equivalently used in \texttt{union} of pointers), permit compilers to
 do arbitrarily sophisticated optimizations.
 
\bigskip


%Additional content of this \S\ref{subsec:analysiot} will be written for the final D1.3\textsuperscript{v2}.


%%%%%%%%%%%%%%%%%%%%%%%%%%%%%%%%%%%%%%%%%%%%%%%%%%%%%%%%%%%%%%%%
%% Local Variables: ;;
%% compile-command: "cd ..; ./build-bismon-doc.sh" ;;
%% End: ;;
%%%%%%%%%%%%%%%%%%%%%%%%%%%%%%%%%%%%%%%%%%%%%%%%%%%%%%%%%%%%%%%%


\newpage
%% file using-bm.tex, which is \input from bismon-chariot-doc.tex
\section{Using \emph{Bismon}}
\label{sec:using}

This section §\ref{sec:using} should become somehow a user manual, and
will be written for the final D1.3\textsuperscript{v2}. It is both for
the ordinary IoT developer just using \emph{bismon} for static
analysis of IoT source code, and for the static analysis expert
configuring and programming it.

Most of that should be generated from data persisted inside \emph{bismon}. Perhaps
should be exchanged with the ``static analysis'' chapter
(§\ref{sec:staticanalys}).

%\bigskip

%Notice that most of \emph{GCC MELT} documentation was generated. See its
%{\small\bmurl{http://starynkevitch.net/Basile/gcc-melt/meltpluginapi.html/meltpluginapi.html}}. Likewise,
%most of this chapter will be machine generated from the persistent
%state of \emph{Bismon}


\subsection{How JSON is used by Bismon}
\label{subsec:json-usage}

The JSON\footnote{See \bmurl{http://json.org/} for more} \index{JSON}
textual format is a convenient, common and compact structured textual
format. It is used in \emph{Bismon}, in particular because of its web
interface, and supported as a payload \index{payload} (but not
directly\footnote{Adding immutable JSON values as a new kind of Bismon
  value could be considered in the future.} as an immutable value) for
objects of class \texttt{json\_object}.  \index{json
  object@\texttt{json\_object}}

Conceptually, the JSON model is close, but not identical to, the
Bismon persistent model: it provides structured and compositional
\index{object}%
\index{attribute}%
\index{component}%
\index{payload}%
constructs, but JSON objects have \emph{strings} as attributes, while
Bismon objects have arbitrary object references as attributes, and
also components and some optional payload.

\subsubsection{The canonical \emph{JSON} encoding of \emph{Bismon} values}
\label{subsubsec:canonical-json-encoding}

Therefore, there is some way to encode any \emph{Bismon} value into a
JSON value; this is the \emph{canonical JSON encoding of
  values}\index{canonical JSON encoding}, given in figure
\ref{fig:canon-json} below.

\begin{figure}[!htbp]
  \begin{relsize}{-1}
  \begin{center}
    \begin{tabular}{lclp{5cm}}
      %%%
      $\llbracket \textrm{nil} \rrbracket_{json}$ \index{nil} &
      $\rightarrow$ &
      \texttt{\textbf{null}}  &
      {\relsize{-1}{The Bismon nil is encoded as the JSON null}} \\
      %%%
      $\llbracket \textrm{unspecified} \rrbracket_{json}$ &
      $\rightarrow$ &
      \texttt{\textbf{false}}  &
      {\relsize{-1}{The Bismon \emph{unspecified} \index{unspecified} is encoded as the JSON false}} \\
      %%%
      $\llbracket \textrm{integer}~ i \rrbracket_{json}$ &
      $\rightarrow$ &
      $i$ {\relsize{-1}{(JSON integer)}} &
      {\relsize{-1}{tagged integers \index{integer} encoded as is}} \\
      %%%
      $\llbracket \textrm{boxed float}~ f \rrbracket_{json}$ &
      $\rightarrow$ &
      $f$ {\relsize{-1}{(JSON float)}} &
      {\relsize{-1}{boxed doubles \index{double} encoded as is, with decimal point}} \\
      %%%
      $\llbracket \textrm{string}~ s \rrbracket_{json}$ &
      $\rightarrow$ &
      $s$  {\relsize{-1}{(JSON string)}}  &
      {\relsize{-1}{Bismon strings \index{string} encoded as is}} \\
      \rule{0pt}{3ex} 
      %%%
      $\llbracket \textrm{object}~ \omega$ {\relsize{-1.5}{\textrm{of objid}~ \textit{oid}}}$ \rrbracket_{json}$ \index{object} \index{objid} &
      $\rightarrow$ &
      \texttt{\textbf{\{ "!oid" :}} $\mathit{oid}$ \texttt{\textbf{\}}} &
             {{\relsize{-1}{Bismon objects encoded with \texttt{"!oid"} JSON attribute giving the \index{objid}objid as a JSON string}}} \\
              \rule{0pt}{3ex}
      %%%
      $\llbracket \textrm{tuple}~ [ \omega_1 \ldots \omega_n ] ~ \rrbracket_{json}$ \index {tuple} &
      $\rightarrow$ & %
      \begin{minipage}[c]{4cm}
         \texttt{\textbf{\{ "!tup" :}} \\
         \hspace*{1.2em}  $\mathtt{\texttt{[}} ~ \mathit{oid}_1 \ldots \mathit{oid}_n ~ \mathtt{\texttt{]}}$ \texttt{\textbf{\}}}
      \end{minipage}
      & {\relsize{-1}{Bismon tuples encoded  with \texttt{"!tup"} JSON attribute giving the JSON array of corresponding objid JSON strings : ~  \mbox{$\mathit{oid}_i = \mathrm{objid} (\omega_i)$}
            }}
      \\
       \rule{0pt}{3ex}
      %%%
      $\llbracket \textrm{set}~ \{ \omega_1 \ldots \omega_n \} ~ \rrbracket_{json}$ \index{set} &
      $\rightarrow$ & %
      \begin{minipage}[c]{4cm}
         \texttt{\textbf{\{ "!set" :}} \\
         \hspace*{1.2em}  $\mathtt{\texttt{[}} ~ \mathit{oid}_1 \ldots \mathit{oid}_n ~ \mathtt{\texttt{]}}$ \texttt{\textbf{\}}}
      \end{minipage}
      & {\relsize{-1}{Bismon sets encoded  with \texttt{"!set"} JSON attribute giving the JSON array of corresponding objid JSON strings : ~  \mbox{$\mathit{oid}_i = \mathrm{objid} (\omega_i)$}
            }}
      \\
       \rule{0pt}{3ex}
      %%%
      $\llbracket \textrm{node}~ \mathtt{\textbf{*}} ~ \omega_{conn} \mathtt{(} \sigma_1 \ldots \sigma_n \mathtt{)} ~ \rrbracket_{json}$ \index{node} &
      $\rightarrow$ & %
      \begin{minipage}[c]{4.5cm}
        \texttt{\textbf{\{ "!node" :}} $\mathit{oid}_{conn}$ \texttt{\textbf{,}} \\
         \hspace*{0.8em}  \texttt{\textbf{"!sons" :}} \\
         \hspace*{1.2em} %
         $\mathtt{\texttt{[}} ~ \llbracket \sigma_1 \rrbracket_{json} ~ \ldots ~  \llbracket \sigma_n \rrbracket_{json} ~ \mathtt{\texttt{]}}$
         \texttt{\textbf{\}}}
      \end{minipage}
      & {\relsize{-1}{Bismon nodes encoded  with \texttt{"!node"} JSON attribute giving the objid $\mathit{oid}_{conn} = \mathrm{objid} (\omega_{conn})$ of the connective $\mathit{oid}_{conn}$, and with  \texttt{"!sons"} JSON attribute associated to the array of encodings of that node's sons $\sigma_i$}
            }
      \\
       \rule{0pt}{3ex}
      %%%
      $\llbracket \textrm{closure}~ \mathtt{\textbf{\%}} ~ \omega_{rout} \mathtt{(} \kappa_1 \ldots \kappa_n \mathtt{)} ~ \rrbracket_{json}$ \index{closure} &
      $\rightarrow$ & %
      \begin{minipage}[c]{4.5cm}
        \texttt{\textbf{\{ "!clos" :}} $\mathit{oid}_{rout}$ \texttt{\textbf{,}} \\
         \hspace*{0.8em}  \texttt{\textbf{"!cval" :}} \\
         \hspace*{1.2em} %
         $\mathtt{\texttt{[}} ~ \llbracket \kappa_1 \rrbracket_{json} ~ \ldots ~  \llbracket \kappa_n \rrbracket_{json} ~ \mathtt{\texttt{]}}$
         \texttt{\textbf{\}}}
      \end{minipage}
      & {\relsize{-1}{Bismon closures encoded  with \texttt{"!clos"} JSON attribute giving the objid  $\mathit{oid}_{rout} = \mathrm{objid} (\omega_{conn})$ of the closure's routine, and with  \texttt{"!cval"} JSON attribute associated to the array of encodings of that closure's closed values $\kappa_i$}
            } \\
      \hline \\
      \rule{0pt}{6ex} \\
      \begin{minipage}[c]{3.5cm}
        $ \nu = $ \\ $\mathrm{apply}  ~ \mathtt{\textbf{\%}} ~ \omega_{rout} \mathtt{(} \kappa_1 \ldots \kappa_n \mathtt{)} $ \\
        to $\mathit{ctxt}, \mathit{depth}$ 
      \end{minipage} &
      $\leftarrow$ & %
      \begin{minipage}[c]{4.5cm}
        \texttt{\textbf{\{ "!apply" :}} $\mathit{oid}_{rout}$ \texttt{\textbf{,}} \\
         \hspace*{0.8em}  \texttt{\textbf{"!cval" :}} \\
         \hspace*{1.2em} %
         $\mathtt{\texttt{[}} ~ \mathit{json}_1 \ldots \mathit{json}_n ~ \mathtt{\texttt{]}}$
         \texttt{\textbf{\}}}
      \end{minipage}
      & {\relsize{-1}{application of object routine $\omega_{rout}$ whose objid is $\mathit{oid}_{rout}$ with closed values $\kappa_1 = \mathit{encode}_{json} ( \mathit{json}_1 )$ \ldots  $\kappa_n = \mathit{encode}_{json} ( \mathit{json}_n )$ to the context, the depth  (during decoding) \textcolor{red}{@@NOT SURE}}}
      \\
       \rule{0pt}{3ex}
    \end{tabular}
    \smallskip
  \end{center}
  \end{relsize}
  \caption{canonical JSON encoding $\llbracket v \rrbracket_{json}$ of a Bismon value $v$.}
  \label{fig:canon-json}
\end{figure}

The canonical JSON encoding is implemented\footnote{Coded in C, in file
\href{https://github.com/bstarynk/bismon/blob/master/jsonjansson\_BM.c}{\texttt{jsonjansson\_BM.c}}}
as the \texttt{canonjsonifyvalue\_BM} \index{canonjsonifyvalue!\texttt{canonjsonifyvalue\_BM}} function.

\medskip

\subsubsection{The nodal \emph{JSON} decoding into \emph{Bismon} values}
\label{subsubsec:nodal-json-decoding}

Since JSON is a structured and compositional, \index{tree!JSON}
tree-like, representation, and because \index{node} nodes are the only
kind of structured immutable \emph{Bismon} values,
\index{decoding!JSON nodal} \index{nodal!JSON decoding} any JSON value
can obviously be decoded into a \emph{Bismon} values, using mostly
nodes for structuring data, following the rules listed in
\ref{fig:nodal-json} below. Actually, there are several variants of
nodal decodings, depending on how JSON strings looking like
\index{objid} full objids (e.g. JSON
\texttt{"\_756o00yB7Zs\_1USbaS25hxl"}), or abbreviated objids
(e.g. JSON \texttt{"€\_9Z2BgJbf4"}), or named objects (e.g.  JSON
\texttt{"arguments"}, related to Bismon object \texttt{arguments},
i.e. \texttt{€\_0jFqaPPHg}) are really decoded.

\begin{figure}[!htbp]
  \begin{center}
    \begin{tabular}{lclp{5.3cm}}
      %%%
      $\langle \mathtt{\textit{\texttt{null}}} \rangle^{nod}$ \index{null} &
      $\rightarrow$ &
      \texttt{\textbf{json\_null}}  &
             {\relsize{-1}{The JSON \texttt{null}\index{null!JSON} is nodal-decoded as the \texttt{json\_null} Bismon object \texttt{€\_6WOSg1mpN}}}
      \\
      %%%
      $\langle \mathtt{\textit{\texttt{false}}} \rangle^{nod}$ \index{null} &
      $\rightarrow$ &
      \texttt{\textbf{json\_false}}  &
             {\relsize{-1}{The JSON \texttt{false} is nodal-decoded as the \texttt{json\_false} Bismon object \texttt{€\_1h1MMlmQi}}}
      \\
      %%%
      $\langle \mathtt{\textit{\texttt{true}}} \rangle^{nod}$ \index{null} &
      $\rightarrow$ &
      \texttt{\textbf{json\_true}}  &
             {\relsize{-1}{The JSON \texttt{true} is nodal-decoded as the \texttt{json\_true} Bismon object \texttt{€\_0ekuRPtKaI}}}
      \\
      %%%
      $\langle \mathrm{integer} ~ \iota \rangle^{nod}$  &
      $\rightarrow$ &
      tagged integer $\iota$  &
             {\relsize{-1}{The JSON integers are nodal-decoded as the corresponding \emph{Bismon} integer}}
      \\
      %%%
      $\langle \mathrm{real} ~ \delta \rangle^{nod}$  &
      $\rightarrow$ &
      boxed double $\delta$  &
             {\relsize{-1}{The JSON reals are nodal-decoded as the corresponding \emph{Bismon} boxed double}}
      \\
      \rule{0pt}{2ex} \\
      %%%
      $\langle \mathrm{\textrm{objid-looking} ~ string ~} \sigma \rangle^{nod}$ \index{objid} &
      $\rightarrow$ &
      \begin{minipage}{4.7cm}
        object $\omega$, or else a \\
        node $\textbf{*}\; \mathtt{\textbf{\texttt{object}}} (\omega) $ with an \\
        object $\omega \in \{ \mathtt{json\_null},$ \\
        \hspace*{0.4em} $ \mathtt{json\_false},$ $  \mathtt{json\_true},$ \\
        \hspace*{0.4em} $\mathtt{json\_array}, \mathtt{json\_object} \}$
      \end{minipage}
      &
      {\relsize{-1}{an objid-looking string $\sigma$ \emph{may} be nodal-decoded as
          \emph{the existing} object $\omega$ such as
          $\mathrm{objid}(\omega) = \sigma$} or else as a node of connective \texttt{object} if $\omega$ is a special object mentioned here.}
      \\
      %%%
      $\langle \mathrm{\textrm{name-looking} ~ string ~} \sigma \rangle^{nod}$ \index{objid} &
      $\rightarrow$ &
      \begin{minipage}{4.7cm}
        object $\omega$, or else a \\
        node $\textbf{*}\; \mathtt{\textbf{\texttt{object}}} (\omega) $ with an \\
        object $\omega \in \{ \mathtt{json\_null},$ \\
        \hspace*{0.4em} $ \mathtt{json\_false},$ $  \mathtt{json\_true},$ \\
        \hspace*{0.4em} $\mathtt{json\_array}, \mathtt{json\_object} \}$
      \end{minipage} &
      {\relsize{-1}{a name-looking string $\sigma$ is nodal-decoded as the \emph{existing named} object $\omega$ such as $\mathrm{name}(\omega) = \sigma$, or else, for names used here, is decoded as
          \emph{a node}  \mbox{$\textbf{*}\; \mathtt{\textbf{\texttt{object}}} ( \omega) $} such as
          $\mathrm{name}(\omega) = \sigma$}}
      \\
      \rule{0pt}{2ex} \\
      %%%
      $\langle \mathrm{\textrm{any} ~ string} ~ \sigma \rangle^{nod}$ \index{objid} &
      $\rightarrow$ &
      \emph{Bismon} string $\sigma$  &
      {\relsize{-1}{a string $\sigma$ \emph{would} otherwise be nodal-decoded as
          is into the same \emph{Bismon string} $\sigma$}}
      \\
      %%%
      $\langle \mathtt{\textbf{[}} \mathit{\,js}_1 \mathtt{\textbf{,}} \ldots \mathit{\,js}_n \mathtt{\textbf{]}}  \rangle^{nod}$ &
      $\rightarrow$ &
      \begin{minipage}{4.7cm}
      $\mathtt{\texttt{\textbf{*}}} \mathtt{\texttt{\textbf{json\_array}}}$ \\
       \hspace*{1.5em} $ ( \langle \mathit{\,js}_1 \rangle^{nod} \ldots  \langle \mathit{\,js}_n \rangle^{nod} ) $
      \end{minipage}
      &
       {\relsize{-1}{A JSON array is compositionally nodal-decoded into a node of connective \texttt{json\_array} and sons given by nodal-decoding the components of that array}} \\
      \rule{0pt}{6ex} \\
      %%%
      $\langle \mathtt{\textbf{\{}} \alpha_1 \mathtt{\textbf{:}} \mathit{\,js}_1 \mathtt{\textbf{,}} \ldots \alpha_n \mathtt{\textbf{:}} \mathit{\,js}_n \mathtt{\textbf{\}}}  \rangle^{nod}$ &
      $\rightarrow$ &
      \begin{minipage}{4.7cm}
      $\mathtt{\texttt{\textbf{*}}} \mathtt{\texttt{\textbf{json\_object}}}$ \\
       \hspace*{1.5em} $ ( \langle \mathit{\,js}_1 \rangle^{nod} \ldots  \langle \mathit{\,js}_n \rangle^{nod} ) $
      \end{minipage}
      &
       {\relsize{-1}{A JSON object is compositionally decoded into a node of connective \texttt{json\_array} and sons given by \texttt{json\_entries} subnodes }} \\
      \rule{0pt}{6ex} \\
    \end{tabular}
  \end{center}
  \caption{nodal JSON decoding $\langle \mathit{js} \rangle^{nod}$ of a JSON value $\mathit{js}$.}
  \label{fig:nodal-json}
\end{figure}

{\color{red}{@@ to be completed a lot,
    explaining conversions of Bismon values to and from JSON.}}

\subsection{Web interface internal design}
\label{subsec:webinterf}

The Web interface \index{web interface} of \emph{bismon} is supposed
to be used without malice (see
§\ref{subsubsec:leveraging-static-analysis} and
§\ref{subsubsec:about-bismon}), with a recent graphical web
browser~\footnote{Such as Firefox 60.5 or later, or Google Chrome 72.0
  or later, both exist in 2019Q1.}  using HTTP/1.1. In particular,
\emph{bismon} does not take any measure against
\index{denial-of-service} denial-of-service attacks, since it is
supposed to be used on a trusted and friendly corporate intranet or
local area network, not directly on the wild Internet. The network
administrator running \emph{bismon} could deploy usual relevant
techniques (firewalls, \texttt{iptables}, HTTP proxying, DMZ, etc ...)
to avoid such attacks. In practice, there are few web browsers - so
few HTTP clients - interacting with \emph{bismon} simultaneously :
only a dozen of people in some IoT development team, and each uses
his/her graphical browser \index{browser} - a \emph{recent} Firefox or
Chrome~\footnote{So no particular effort is even taken to support a
  variety of old browsers: we don't have any code to e.g. support
  Internet Explorer pecularities or deficiencies. Likewise,
  scalability to thousands of simultaneous HTTP connections is out of
  scope in \emph{bismon}, but it is essential in most web
  applications.}. Each \emph{Bismon} user is expected to have one, or
only a few, browser tab[s] interacting with the \texttt{bismon}
server, and these tabs, if there are more than one, are handled as
different web browsers so have different web sessions.  They are
physically and geographically located on the same local area network
as the machine running the \emph{bismon} monitor. So, from web
technologies perspective, \emph{bismon} is making different
trade-offs~\footnote{For example, we could accept making some HTTP
  exchange - e.g. with AJAX - on \emph{every} keystroke on the
  keyboard, but such practice won't be acceptable in usual web
  services. Also, we don't care much about minimizing the HTTP
  exchanges - no ``minification'' needed in practice!} than
``traditional'' web servers or web applications : the web browser
$\leftrightarrow$ \emph{bismon} web server round-trip transmission
time is supposed to be very small so frequent AJAX requests are
possible, the bandwidth is expected to be quite large so voluminous
HTTP responses are acceptable, the number of simultaneous web
connections or of web sessions is tiny. Therefore most web
optimizations are practically unneeded.

With its initial (and current, in mid-2019) naive stop-the-world
garbage collector, the interactive performance and user experience
(i.e. user look-and-feel) of \emph{Bismon} is expected to be
unsatisfactory (since that GC could ``block'' the \texttt{bismon}
monitor and web service for more than half a second - during which the
web interface stays unresponsive, if running the GC on a large enough
heap; but see footnote \ref{fn:initial-gc} suggesting an
improvement). With significant work, that could be improved.

Each HTTP request \index{request} \index{web request} either
corresponds to a ``static'' file path under \texttt{webroot/} (for a
\texttt{GET} or \texttt{HEAD} HTTP request) or else it is handled
dynamically. For a static file path, that file is served directly by
\texttt{onion\_handler\_export\_local\_new} with a
\texttt{Content-Type} corresponding to its suffix; for example an HTTP
\texttt{GET} request of \texttt{favicon.ico} is answered with the
content of \texttt{webroot/favicon.ico} file, and an HTTP \texttt{GET}
request of \texttt{jscript/jquery.js} is served by the content of
\texttt{webroot/jscript/jscript.js}. Care is taken~\footnote{In
  particular, any HTTP request containing \texttt{..} is rejected.} to
avoid serving any static file outside of \texttt{webroot/}. So the
\texttt{webroot/} directory contains static content such as images,
external JavaScript libraries, CSS stylesheets, etc... Static content
requests are always handled the same, so they work even without any
cookies. \index{request!static} \index{static request}

Any HTTP request which cannot be handled as a static resource like
above, because it has no corresponding file under \texttt{webroot/},
is considered as a request for dynamic content and is called a
\index{request!dynamic} \index{dynamic request} \emph{dynamic
  request}. Dynamic content requires a web session \index{web session}
\index{session!web session} cookie \index{web cookie} \index{cookie}
named \texttt{BISMONCOOKIE} which contains~\footnote{A practical
  example of \texttt{BISMONCOOKIE} value might be
  \texttt{n000041R970099188t330716425o\_6IHYL1fOROi\_58xJPnBLCTe}: 41
  is the serial number counting web sessions in the running
  \emph{bismon} process, 970099188 and 330716425 are two random
  numbers, \texttt{\_6IHYL1fOROi\_58xJPnBLCTe} is the
  randomly-generated objid of the web session object.} a cryptographic
quality hash (in practice unforgeable) and mentions the objid
\index{objid} (cf §\ref{subsubsec:objects}) of some web session
object. If there is no cookie, or if the \index{cookie} cookie is
invalid or wrong (e.g. forged), a login \index{login} form is
returned. So any HTTP request for a dynamic content (that is which is
not handled as a static resource like above) is rejected (with HTTP
status \texttt{403 Forbidden}) if the user (a \index{contributor}
contributor in \emph{bismon} parlance,
cf. §\ref{subsubsec:bismon-evolving}) is not logged in.

Dynamic requests are reified as very temporary \emph{bismon} objects
of class \texttt{webexchange\_object}. Their \emph{web exchange}
payload~\footnote{There is no programmatic way to create such a web
  exchange payload. It can only be created by processing such dynamic
  HTTP requests.} contains not only a string buffer~\footnote{Since a
  string buffer should contain valid UTF-8 string content without nul
  bytes, this restriction forbids binary contents in HTTP replies to
  dynamic requests. Hence, dynamically computed image contents are not
  possible, unless they use a textual format like SVG.}, to be filled
with the HTTP response content, but also mentions the web request (as
processed by \texttt{libonion}) and the web session object computed
from the \texttt{BISMONCOOKIE} and may contain some arbitrary data
value. The web exchange object is supposed to be filled -like string
buffers are- and at last given some integer HTTP status and
immediately sent back to the browser.  Their web session object is
created at web login time and is of class
\texttt{websession\_object}~\footnote{So the only way to create a web
  session payload is thru the login form. There is no programmatic way
  to create it.}. It knows the contributor who is logged in, the
expiration time of the session, some session data (an arbitrary
\emph{bismon} value; of course more data can sit in attributes or in
components of that web session object), and the web socket \index{web
  socket} connection (if any) to the browser using that session. The
session storage~\footnote{See
  \url{https://developer.mozilla.org/en-US/docs/Web/API/Window/sessionStorage}
  for more.} associated to key \texttt{"bismontab"} identifies and
gives the tab number in the browser. An inactive web session expires
in about an hour~\footnote{See the
  \texttt{USER\_WEBSESSION\_EXPIRATION\_DELAY} constant in
  \texttt{web\_ONIONBM.c}}.

Of course, web request objects or web session objects are
\index{transient} transient and are not and should not be persisted at
dump \index{dump} time (cf. §\ref{subsec:persistence}). So after each
restart of the \emph{bismon} monitor, its web users
(i.e. contributors) should login \index{login} again.

A dynamic request is handled by some closure \index{closure} and
should be answered in a couple~\footnote{See the
  \texttt{WEBEXCHANGE\_DELAY\_BM} constant in file
  \texttt{web\_ONIONBM.c} ...} of seconds; otherwise a \index{web
  timeout} web timeout \index{timeout!web} occurs. That web handler
\index{web handler} closure is applied to the remaining URL path
string and to the web exchange object created in the
\texttt{libonion}-specific thread dealing with the HTTP request, so
outside of the agenda machinery (cf §\ref{subsec:multithreadist}), and
usually would add some \index{tasklet} tasklet into the \index{agenda}
agenda. Most of the time, a fraction of a second later, some other
tasklet would complete the filling the web request object and give
some HTTP status code such as \texttt{200 OK}, then an HTTP reply is
sent back to the browser. If a timeout occurs because the web request
object has not been taken care of quickly enough, an HTTP \texttt{500
  Internal Server Error} is given back to the browser and that web
request object is cleared.

The mapping between URL paths (or prefixes) and web handler closures
handling dynamic requests for them is given by the
\texttt{webdict\_root}~\footnote{For example: an URL like
  \texttt{http://localhost:8086/show/status} is handled by some
  \emph{bismon} monitor listening HTTP requests on port 8086 and with
  \texttt{webdict\_root} associating the string \texttt{"show"} to
  some closure $\kappa$, that web handler closure $\kappa$ would be
  applied to the suffix string \texttt{"status"} and to the web
  exchange object $\omega$ created for that HTTP request. The result
  of that application is ignored, only side effects -often adding some
  tasklets into the agenda, and/or filling the web exchange object
  with some XHTML5, etc...- are useful. If the string in such a web
  handling dictionnary is associated to some other object
  $\omega_{wh}$ of class \texttt{webhandler\_dict\_object}, that
  dictionnary object $\omega_{wh}$ is recursively explored with the
  rest of that URL path (e.g. \texttt{"status"} in our example).}
dictionnary predefined object; for an empty path in the URL (such as
{\relsize{-1}{\texttt{http://localhost:8086/}}} for example), its
\texttt{web\_empty\_handler} attribute is~\footnote{Since dictionnary
  objects map \emph{non-empty} strings to non-nil values
  (cf. §\ref{subsubsec:objects}).} used.  If finally no web handler
closure is found, an \texttt{404 Not found} status is returned. The
\texttt{the\_web\_sessions} predefined object stores the dictionnary
of transient web session objects and associates a cookie string to its
web session object. That dictionnary is forcibly cleared at start of
the web server inside \emph{bismon}, but it should be loaded empty,
since web session objects are created and should remain transient.

In practice, dynamic requests are usually generating the HTML5
\index{HTML5} content very dynamically. For generated HTML, it is much
easier to produce \index{XHTML5} XHTML5, the XML variant \index{XML}
of HTML5, because its textual syntax is~\footnote{For instance, within
  a \texttt{<script} HTML5 element containing JavaScript, it is not
  even allowed in HTML5 to have \texttt{if (x \&lt; 5)} even if
  ordinary HTML rules suggest to use \texttt{\&lt;} instead of
  \texttt{<} in textual content... That makes compositional generation
  of mixture of HTML and JavaScript emitting HTML much harder.} much
more regular and easier to generate than with plain HTML5.

The \texttt{webxhtml\_module} \index{module!\texttt{webxhtml\_module}}
in \emph{bismon} has code to ease the emission of XHTML5. And XHTML5
fragments are emitted by the \texttt{emit\_xhtml} routine
\index{emit xhtml@\texttt{emit\_xhtml}}
object~\footnote{So that \texttt{emit\_xhtml} is, like PHP, a
  machinery to emit arbitrary XHTML. However, we want to avoid
  thinking -like PHP was originally designed- in terms of emitting a
  stream of characters, and \texttt{emit\_xhtml} is supposed to emit
  structured XHTML from some structured, tree-like, internal
  representation. That internal representation is a DAG (directed
  acyclic graph).}. That \texttt{emit\_xhtml}, which get as arguments:
the value \texttt{v\_html} to emit; an arbitrary web context object
\texttt{o\_emitctx} which might be in simple cases just some web
session object; a string buffer object \texttt{o\_strbuf} which is
often the web exchange object; the tagged integer recursion depth
\texttt{v\_depth}, which is in general \emph{not}~\footnote{In very
  simple cases, without closures or sequences in the DAG of emitted
  values, the depth could be the depth of XHTML elements, so could be
  the indentation. In general, it is not.} the emitted
  indentation. When the string buffer is too full or the recursion
  depth is too deep, that \texttt{emit\_xhtml} fails. When the emitted
  HTML-reifying value \texttt{v\_html} is nil, nothing happens. When
  it is a scalar, it is emitted trivially: a string is emitted
  HTML-encoded (so \texttt{\&lt;} for \texttt{<}, etc...); a tagged
  integer is emitted in decimal notation; When \texttt{v\_html} is an
  object $\omega_{html}$, it is emitted per the following rules:

  \begin{itemize}
    
\item the \texttt{newline} object emits an indented newline.
  \index{newline@\texttt{newline}}
  
\item the \texttt{nlsp} object emits a newline when the current line
  is long enough, or else a space.
  \index{nlsp@\texttt{nlsp}}
  
\item the \texttt{space} object emits a space character.
  \index{space@\texttt{space}}
  
\item instances of \texttt{html\_void\_element\_object} emit some void
\index{html void element object@\texttt{html\_void\_element\_object}}
  element like e.g. \texttt{<hr class='foo'/>} using the
  \texttt{€\_0FRLxSGQlZ} routine. The \texttt{emit\_xhtml\_open}
\index{emit xhtml open@\texttt{emit\_xhtml\_open}}
  selector should emit -as a side effect- the opening tag \texttt{<hr
    class='foo'} without the ending \texttt{/>}, and returns the string naming the tag, e.g. \texttt{"hr"}.
  
\item instances of \texttt{html\_element\_object} emit recursively
\index{html element object@\texttt{html\_element\_object}}
  using \texttt{€\_5NH940lCHYJ} some nested XHTML element starting
  with a start tag like \texttt{<div} but ending with an end tag like
  \texttt{</div>}; the components of that objects are emitted
  recursively (with an incremented recursion depth). The spacing style
\index{html spacing@\texttt{html\_spacing}}
  is first determined by sending \texttt{html\_spacing} with the
  \texttt{o\_emitctx} and the depth. It can be \texttt{newline} for
  indented, newline separated, content; or \texttt{nlsp} for space or
  newline separated, unindented content; or \texttt{space} for space
  separated content; any other spacing style -notably nil- don't emit
  any separators in the content. The start tag is emitted with
  \texttt{emit\_xhtml\_open} returning the tag string like before
  (e.g. \texttt{"span"} for a \texttt{<span class='somecl'} emission,
  then \texttt{>} is output to end the opening tag, then the
  components, then the end tag is emitted, using the returned tag
  string from \texttt{emit\_xhtml\_open}.
  
\item instances of \texttt{html\_sequence\_object} emit
      recursively their components but without surrounding tags.
      \index{html sequence object@\texttt{html\_sequence\_object}}
      
\item instances of \texttt{html\_active\_object} emit recursively HTML
      \index{html active object@\texttt{html\_active\_object}}
      stuff thru a message of selector \texttt{emit\_xhtml} sent to them.
      
\item any other object is emitted by its name, if it has some, or by
  its objid. This is mostly intended to represent common repeated names or
  words by a single and shared object.
\end{itemize}

When  \texttt{v\_html} is a node of connective $\omega_{conn}$, it is emitted per the following rules:

\begin{itemize}

  \index{int@\texttt{int}}
\item if $\omega_{conn}$ is one of \texttt{int}, \texttt{hexa},
  \index{int@\texttt{int}}
  \index{hexa@\texttt{hexa}}
  \index{octa@\texttt{octa}}
  \texttt{octa} and \texttt{v\_html} is an unary node with a integer
  son $n$, that integer $n$ is emitted in decimal, hexadecimal, octal respectively.

\item if $\omega_{conn}$ is \texttt{id}~\footnote{So nodes of
  connective \texttt{id}, \texttt{object} or \texttt{name} can be used
  to emit objects of class \texttt{html\_void\_element\_object},
  \texttt{html\_element\_object}, \texttt{html\_sequence\_object},
  \texttt{html\_active\_object} which would be handled specially
  otherwise.}  and \texttt{v\_html} is an unary node with an object
  \index{id@\texttt{id}}
  \index{object@\texttt{object}}
  \index{name@\texttt{name}}
  son $\omega_{son}$ its objid is emitted.
  
\item if $\omega_{conn}$ is \texttt{buffer} and \texttt{v\_html} is an
  \index{buffer@\texttt{buffer}}
  unary node with an object son $\omega_{son}$ which has a string
  buffer payload, that string is emitted HTML-encoded.
  
\item if $\omega_{conn}$ is \texttt{object} and \texttt{v\_html} is an unary node with an object son $\omega_{son}$ its name or objid is emitted.
  \index{object@\texttt{object}}
  
\item if $\omega_{conn}$ is \texttt{name} and \texttt{v\_html} is an
  unary node with an object son $\omega_{son}$ its name is emitted,
  \index{name@\texttt{name}}
  and when $\omega_{son}$ is not a named object, we have a failure. If $\omega_{conn}$ is \texttt{name} and \texttt{v\_html} is an
  binary node with an object first son $\omega_{son}$, and some arbirary non-nil second son $\epsilon$ its name is emitted,
  and when $\omega_{son}$ is not a named object, the $\epsilon$ is recursively emitted
  
\item if $\omega_{conn}$ is \texttt{sequence}, every son of
  \texttt{v\_html} is emitted in sequence, with an incremented
  $depth$. Nothing is additionally emitted between
  them. \index{sequence@\texttt{sequence}}

\item if $\omega_{conn}$ is \texttt{space}, or \texttt{newline}, or
  \texttt{nlsp}, every son of \texttt{v\_html} is emitted in sequence,
  with an incremented $depth$. Between each son, a space
  (respectively, a newline, or a smart space of newline) is emitted.
  \index{space@\texttt{space}} \index{newline@\texttt{newline}}
  \index{nlsp@\texttt{nlsp}}


\item for any other object $\omega_{conn}$ as connective, we extract
  its \texttt{emit\_xhtml\_node} attribute $v_{emit~node}$ and its
  \texttt{emit\_xhtml\_connective\_open} attribute
  $v_{emit~open}$. Only one of them should be present (non-nil value)
  and it should be a closure. If $v_{emit~node}$ is given, it is
  applied to $\omega_{conn}$ \texttt{o\_emitctx o\_strbuf $depth + 1$
    v\_html}, else if $v_{emit~open}$ is present, we apply it (like
  the \texttt{emit\_xhtml\_open} selector above) to $\omega_{conn}$
  \texttt{o\_emitctx o\_strbuf $depth + 1$ v\_html} to obtain an XHTML
  element tag. We also extract and use its \texttt{html\_spacing}
  attribute. Then proceed like for \texttt{html\_element\_object}
  using the sons as components....
  
\item {\color{red}{@@ to be completed a lot.}}
  
\end{itemize}


When \texttt{v\_html} is a closure, it is applied {\color{red}{@@ to
    be completed}} and the result of that application is recursively
emitted. When \texttt{v\_html} is a sequence (set or tuple), its
components are emitted recursively.

If no rule is applicable, \texttt{emit\_xhtml} fails.

\medskip

The web session objects are also used for \index{WebSocket}
\emph{WebSockets} with the following additional conventions. The
\texttt{bismon} server uses WebSockets only for asynchronous
communication from that \texttt{bismon} server to Web
browsers~\footnote{So web browsers don't communicate
  \emph{asynchronously} with the \texttt{bismon} server. For such
  communications from browser to \texttt{bismon}, Web browser always
  use synchronous HTTP requests, e.g. using AJAX techniques.}. The
WebSocket messages from \texttt{bismon} to web browsers are arbitrary
\emph{JSON} values.

\subsection{Using \texttt{bismon} for \textsc{Chariot}}
\label{subsec:bismon-for-chariot}

To run the \texttt{bismon} monitor for \textsc{Chariot} related
activities, that monitor should initialize its state for these
activities. So you need to pass \texttt{-i init\_chariotdemo} as a
program argument when running \texttt{bismon} in that case.

%%%%%%%%%%%%%%%%%%%%%%%%%%%%%%%%%%%%%%%%%%%%%%%%%%%%%%%%%%%%%%%%
%% Local Variables: ;;
%% compile-command: "cd ..; ./build-bismon-doc.sh" ;;
%% End: ;;
%%%%%%%%%%%%%%%%%%%%%%%%%%%%%%%%%%%%%%%%%%%%%%%%%%%%%%%%%%%%%%%%


\newpage
% file miscwork-bm.tex, which is \input from bismon-doc.tex
\section{miscellanous work}
\label{sec:miscwork}

\subsection{contributions to other free software projects}

\label{subsec:contribfree}
This is related to subtask ST1.3.2 of \textsc{Chariot} GA.

\subsubsection{aborted contribution to \texttt{libonion}}
\label{subsubsec:contriblibonion}
The \texttt{libonion} library is a free software HTTP server library
(LGPLv3 licensed) that is used in \texttt{bismon} for its web service
feature. See its web site \bmurl{https://www.coralbits.com/libonion/}
for a description, and its source repository
\bmurl{https://github.com/davidmoreno/onion} for more.

The handling of \texttt{SIGTERM} signal (and others) is deemed
unsatifactory. See the opened issue 229 {\relsize{-1}{(on
    \bmurl{https://github.com/davidmoreno/onion/issues/229})}} in
\texttt{libonion}.  We discussed that issue on google group with the
\emph{libonion} community, and came to a disagreement (our design was
considered too complex, but we believe that such a complexity is
needed to avoid bugs in the rare cases of a multi ``\texttt{onion}''
application, which \texttt{bismon} is not).

Independently of that issue, we improved our \texttt{bismon} to avoid
needing or depending on that \texttt{SIGTERM} feature in
\texttt{libonion} (by using \texttt{signalfd} Linux specific
facilities in \texttt{bismon} itself and passing the
\texttt{O\_NO\_SIGTERM} flag to \texttt{onion\_new}...).

So the effort on improving \texttt{SIGTERM} handling in
\texttt{libonion} was concluded.

\subsubsection{contribution to \emph{GCC}}
\label{subsubsec:contribgcc}
There is no contribution yet to \emph{GCC}, because it is not yet
needed in october 2018. We reserve some effort for future such
contributions, when our \emph{GCC} plugin generator would require
them. In the lucky case where no adaptation of \emph{GCC} plugin
infrastructure is necessary, the effort could be moved to other work
in T1.3 (notably ST1.3.3).

\subsection{design and implementation of the compiler and linker extension}

\label{subsec:compilinkext}
This is related to subtask ST1.3.4 (and also ST1.3.1) of \textsc{Chariot} GA 

The compiler extensions will be \emph{generated} GCC plugins.

The linker extension will compute some ``cryptographic quality'' hash
code of the C or C++ translation units of the IoT software. Then it
will interact with the blockchain, according to the \emph{§6 API for
  Private key related transactions} of the \emph{D1.2 Method for
  coupling preprogrammed private keys on IoT devices with a Blockchain
  system}. That API is a Web API and a C or C++ compatible plain API
or library should be developed, following the tutorial code example of
D1.2.

@@TO BE COMPLETED (in D1.3v2).


\newpage
% file conclus-bm.tex, which is \input from bismon-chariot-doc.tex
\section{Conclusion}
\label{sec:conclusion}

The \texttt{bismon} free software is developed in an agile and
incremental manner~\footnote{So there are no released stable versions
  of this software, but snapshots.} (required by its bootstrapping
approach), with continuous updates to
\bmurl{https://github.com/bstarynk/bismon/}.

In october 2018, the persistence machinery is working and daily used
to enhance \texttt{bismon}. The agenda mechanism is working. A naive
stop-the-world mark-and-sweep precise garbage collector is
implemented. The generation of internal C code is done (by
hand-written routines, still coded in C), this enables the
meta-programming approach. The web interface is worked upon: a
\texttt{libonion} based infrastructure is already handling HTTP
requests, and a GDPR-compliant login form is presented on web
browsers. Our \texttt{jsweb\_module} contains the functions related to
Javascript (nearly complete) and HTML generation (work in progress). The
syntactical editor (replacing the crude GTK interface) and then the
GCC plugin generation should be worked on.

In august 2019, the web machinery is mostly working. More generated C
code is available. The JSON handling is incomplete. Bismon
continuations \index{contination!reification} \index{reification!of
  continuations} are almost\footnote{Thanks to generated invocations of
  the \texttt{LOCALFRAME\_BM}
  \index{LOCALFRAME\_BM@\texttt{LOCALFRAME\_BM}} \emph{C} macro, which
  provides 90\% of the development work: full transient reification of
  partial continuations, that is of \index{call stack} \index{call
    frame} call stack segments, is just a matter of clustering emitted
  stack-local \index{stackframe\_stBM@\texttt{stackframe\_stBM}}
  \texttt{struct stackframe\_stBM}-based linked-lists of Bismon call
  frames.}  reifiable into transient \index{transient!object for
  continuations} objects, having as payload a linked-list sequence of
\index{call frame} call frames.

The final D1.3~\textsuperscript{v2} version (scheduled for M30) of
this deliverable will explain the Web interfaces (both for the
ordinary user, i.e. the IoT developer; and for the static analysis
expert) and the generation of C++ code for GCC plugins, with some
examples of simple, IoT focused, whole-program static source code
analysis performed by \emph{bismon}. So the final
D1.3~\textsuperscript{v2} document will contain a longer conclusion.

Machine learning techniques inspired by \cite{zhang:2019:learned}
could be relevant in \textsc{Bismon}. See also the
\href{http://refpersys.org/}{\textsc{RefPerSys}}
research\footnote{Also related:
\href{https://afia.asso.fr/journee-hommage-j-pitrat/}{talks in the
  memory of J.Pitrat, AFIA, March 6\textsuperscript{th}, 2020,
  Paris.}} project.

\medskip



%%%%%%%%%%%%%%%%%%%%%%%%%%%%%%%%%%%%%%%%%%%%%%%%%%%%%%%%%%%%%%%%

\clearpage

%% https://en.wikibooks.org/wiki/LaTeX/Indexing#Adding_index_to_table_of_contents
\addcontentsline{toc}{section}{Index}

\medskip

%\section*{sec Index}

\printindex

\medskip

\clearpage
\addcontentsline{toc}{section}{References}

%\section*{sec References}

\medskip

\bibliography{bismon-biblio}

\begin{flushright}
  \begin{relsize}{-2}
    For books in French, I have provided a \emph{tentative} translation into
    English of their title in brackets.
  \end{relsize}
\end{flushright}

\medskip

\bigskip

\hrule

\bigskip

\subsubsection*{About this document}

\begin{center}
\Ovalbox{
\begin{minipage}{0.9\textwidth}
  \medskip
  \begin{quote}
    % TODO: remove the grey when HeVeA can process this
    %%%%%%%%%%%%%%%%%%%%%%%%%%
    To produce this document, both in PDF {\textcolor{Grey}{and HTML}}
    forms : build \emph{bismon}\footnote{See the
      \href{https://github.com/bstarynk/bismon/blob/master/README.md}{\texttt{README.md}}
      file on \bmurl{https://github.com/bstarynk/bismon/} for building
      instructions.} on your Linux computer, then run \texttt{make
      doc}\footnote{That uses {\LaTeX} and
      \href{http://hevea.inria.fr/}{\emph{HeVeA}}. {\textcolor{Grey}{HTML generation might not work in summer 2019.}} }  or just
    \texttt{make latexdoc} to get only its PDF form.

Feedback and improvements on this document can be suggested by email
(to \bmemail{basile@starynkevitch.net} or
\bmemail{basile.starynkevitch@cea.fr}) or by submitting patches to
\textit{Bismon} thru its \bmurl{https://github.com/bstarynk/bismon}
repository (or directly by email, with your permission to include
it). Notice that this document may contain generated documentation,
and will contain more and more generated parts in the future.
  \end{quote}
\bigskip
\end{minipage}
}
\end{center}

\bigskip

\hrule

\bigskip

\section*{Acknowledgements}

Thanks to my colleague Franck Védrine, to several members of the \textsc{Chariot} consortium who have
proofread this report, and to Niklas Rosencratz (from Sweden, outside
of the consortium) who voluntarily found mistakes in it and proposed,
in the repository on \bmurl{https://github.com/montao/bismon-docker/},
a \texttt{Dockerfile} for \emph{bismon}.

Thanks also to Jérémie Salvucci (France) and Abhishek Chakravarti
(India) for many valuable questions, suggestions, and discussions -in
numerous audioconference (or face to face) exchanges, or in private
emails-, about reflexive persistent systems in general, and more
specifically about \emph{Bismon}. Both did suggest several
improvements to this report and to the \emph{Bismon} software.

\medskip

\hrule 

\medskip

\vspace{2cm}

\bigskip

The CC-BY-SA license below for this (\textsc{Chariot}
D1.3\textsuperscript{v2}) deliverable is required to enable major parts
of this report to be later incorporated into a proper
\emph{bismon} free software documentation. For example, the Debian
Linux distribution has a policy~\footnote{See Debian documentation
  project on \bmurl{https://www.debian.org/doc/ddp} etc...} strongly
recommending a specific set of licenses (notably CC-BY-SA) for
documentation. Using other (deemed proprietary) licenses in free
software documentation is decreasing the future chances of such
documentation being later incorporated in Linux distributions.

\begin{quote}
   \bmincludewidthgraphics{75pt}{CC-BY-SA-4-img}{png}{png} \\
    This entire document
is licensed under the Creative Commons Attribution-ShareAlike 4.0
International License. To view a copy of this license, visit
\bmurl{http://creativecommons.org/licenses/by-sa/4.0/} or send a letter to
Creative Commons, PO Box 1866, Mountain View, CA 94042, USA.
\end{quote}

\bigskip

\vspace{2cm}

\bigskip

\hrule 

\medskip

\begin{flushright}
  \begin{relsize}{-1}
    (git commit \texttt{\bmgitcommit}; generated on
    \textit{\bmdoctimestamp})\\
    Some variant of this draft report is
    downloadable from \bmurl{http://starynkevitch.net/Basile/bismon-chariot-doc.pdf} and
    elsewhere.
  \end{relsize}
\end{flushright}

\end{document}
%%%%%%%%%%%%%%%%%%%%%%%%%%%%%%%%%%%%%%%%%%%%%%%%%%%%%%%%%%%%%%%%
%% Local Variables: ;;
%% compile-command: "cd ..; ./build-bismon-doc.sh" ;;
%% End: ;;
%%%%%%%%%%%%%%%%%%%%%%%%%%%%%%%%%%%%%%%%%%%%%%%%%%%%%%%%%%%%%%%%
